% Options for packages loaded elsewhere
\PassOptionsToPackage{unicode}{hyperref}
\PassOptionsToPackage{hyphens}{url}
\documentclass[
]{article}
\usepackage{xcolor}
\usepackage{amsmath,amssymb}
\setcounter{secnumdepth}{-\maxdimen} % remove section numbering
\usepackage{iftex}
\ifPDFTeX
  \usepackage[T1]{fontenc}
  \usepackage[utf8]{inputenc}
  \usepackage{textcomp} % provide euro and other symbols
\else % if luatex or xetex
  \usepackage{unicode-math} % this also loads fontspec
  \defaultfontfeatures{Scale=MatchLowercase}
  \defaultfontfeatures[\rmfamily]{Ligatures=TeX,Scale=1}
\fi
\usepackage{lmodern}
\ifPDFTeX\else
  % xetex/luatex font selection
\fi
% Use upquote if available, for straight quotes in verbatim environments
\IfFileExists{upquote.sty}{\usepackage{upquote}}{}
\IfFileExists{microtype.sty}{% use microtype if available
  \usepackage[]{microtype}
  \UseMicrotypeSet[protrusion]{basicmath} % disable protrusion for tt fonts
}{}
\makeatletter
\@ifundefined{KOMAClassName}{% if non-KOMA class
  \IfFileExists{parskip.sty}{%
    \usepackage{parskip}
  }{% else
    \setlength{\parindent}{0pt}
    \setlength{\parskip}{6pt plus 2pt minus 1pt}}
}{% if KOMA class
  \KOMAoptions{parskip=half}}
\makeatother
\usepackage{color}
\usepackage{fancyvrb}
\newcommand{\VerbBar}{|}
\newcommand{\VERB}{\Verb[commandchars=\\\{\}]}
\DefineVerbatimEnvironment{Highlighting}{Verbatim}{commandchars=\\\{\}}
% Add ',fontsize=\small' for more characters per line
\newenvironment{Shaded}{}{}
\newcommand{\AlertTok}[1]{\textcolor[rgb]{1.00,0.00,0.00}{\textbf{#1}}}
\newcommand{\AnnotationTok}[1]{\textcolor[rgb]{0.38,0.63,0.69}{\textbf{\textit{#1}}}}
\newcommand{\AttributeTok}[1]{\textcolor[rgb]{0.49,0.56,0.16}{#1}}
\newcommand{\BaseNTok}[1]{\textcolor[rgb]{0.25,0.63,0.44}{#1}}
\newcommand{\BuiltInTok}[1]{\textcolor[rgb]{0.00,0.50,0.00}{#1}}
\newcommand{\CharTok}[1]{\textcolor[rgb]{0.25,0.44,0.63}{#1}}
\newcommand{\CommentTok}[1]{\textcolor[rgb]{0.38,0.63,0.69}{\textit{#1}}}
\newcommand{\CommentVarTok}[1]{\textcolor[rgb]{0.38,0.63,0.69}{\textbf{\textit{#1}}}}
\newcommand{\ConstantTok}[1]{\textcolor[rgb]{0.53,0.00,0.00}{#1}}
\newcommand{\ControlFlowTok}[1]{\textcolor[rgb]{0.00,0.44,0.13}{\textbf{#1}}}
\newcommand{\DataTypeTok}[1]{\textcolor[rgb]{0.56,0.13,0.00}{#1}}
\newcommand{\DecValTok}[1]{\textcolor[rgb]{0.25,0.63,0.44}{#1}}
\newcommand{\DocumentationTok}[1]{\textcolor[rgb]{0.73,0.13,0.13}{\textit{#1}}}
\newcommand{\ErrorTok}[1]{\textcolor[rgb]{1.00,0.00,0.00}{\textbf{#1}}}
\newcommand{\ExtensionTok}[1]{#1}
\newcommand{\FloatTok}[1]{\textcolor[rgb]{0.25,0.63,0.44}{#1}}
\newcommand{\FunctionTok}[1]{\textcolor[rgb]{0.02,0.16,0.49}{#1}}
\newcommand{\ImportTok}[1]{\textcolor[rgb]{0.00,0.50,0.00}{\textbf{#1}}}
\newcommand{\InformationTok}[1]{\textcolor[rgb]{0.38,0.63,0.69}{\textbf{\textit{#1}}}}
\newcommand{\KeywordTok}[1]{\textcolor[rgb]{0.00,0.44,0.13}{\textbf{#1}}}
\newcommand{\NormalTok}[1]{#1}
\newcommand{\OperatorTok}[1]{\textcolor[rgb]{0.40,0.40,0.40}{#1}}
\newcommand{\OtherTok}[1]{\textcolor[rgb]{0.00,0.44,0.13}{#1}}
\newcommand{\PreprocessorTok}[1]{\textcolor[rgb]{0.74,0.48,0.00}{#1}}
\newcommand{\RegionMarkerTok}[1]{#1}
\newcommand{\SpecialCharTok}[1]{\textcolor[rgb]{0.25,0.44,0.63}{#1}}
\newcommand{\SpecialStringTok}[1]{\textcolor[rgb]{0.73,0.40,0.53}{#1}}
\newcommand{\StringTok}[1]{\textcolor[rgb]{0.25,0.44,0.63}{#1}}
\newcommand{\VariableTok}[1]{\textcolor[rgb]{0.10,0.09,0.49}{#1}}
\newcommand{\VerbatimStringTok}[1]{\textcolor[rgb]{0.25,0.44,0.63}{#1}}
\newcommand{\WarningTok}[1]{\textcolor[rgb]{0.38,0.63,0.69}{\textbf{\textit{#1}}}}
\usepackage{longtable,booktabs,array}
\usepackage{caption}
\captionsetup[table]{skip=6pt}
\newcounter{none} % for unnumbered tables
\usepackage{calc} % for calculating minipage widths
% Correct order of tables after \paragraph or \subparagraph
\usepackage{etoolbox}
\makeatletter
\patchcmd\longtable{\par}{\if@noskipsec\mbox{}\fi\par}{}{}
\makeatother
% Allow footnotes in longtable head/foot
\IfFileExists{footnotehyper.sty}{\usepackage{footnotehyper}}{\usepackage{footnote}}
\makesavenoteenv{longtable}
\setlength{\emergencystretch}{3em} % prevent overfull lines
\providecommand{\tightlist}{%
  \setlength{\itemsep}{0pt}\setlength{\parskip}{0pt}}
\usepackage{bookmark}
\IfFileExists{xurl.sty}{\usepackage{xurl}}{} % add URL line breaks if available
\urlstyle{same}
% fallback for those not using the hyperref driver hyperxmp:
\makeatletter
\@ifundefined{xmpquote}{\newcommand{\xmpquote}[1]{#1}}{}
\makeatother
\hypersetup{
  hidelinks,
  pdfcreator={LaTeX via pandoc}}

\author{}
\date{}

\begin{document}

\section{BAB I: METODOLOGI ANALISIS EKSPANSI INDUSTRI EKSTRAKTIF DAN
INFRASTRUKTUR PENUNJANG DI PULAU
SULAWESI}\label{bab-i-metodologi-analisis-ekspansi-industri-ekstraktif-dan-infrastruktur-penunjang-di-pulau-sulawesi}

Dokumen laporan metodologi ini menyajikan kerangka ilmiah, landasan
regulasi, formulasi matematis, prosedur analisis statistik, serta
metodologi pembuktian berbasis data terbuka yang dioperasionalkan pada
\textbf{Bab 1: Ekspansi Industri Ekstraktif} dalam studi Daya Dukung dan
Daya Tampung Lingkungan Hidup (D3TLH) Sulawesi periode 2014--2024.

\begin{center}\rule{0.5\linewidth}{0.5pt}\end{center}

\subsection{SUB-BAB 1.1: Konteks Makro: Breakdown PDRB per
Komoditas}\label{sub-bab-1.1-konteks-makro-breakdown-pdrb-per-komoditas}

\subsubsection{1.1.1 Konteks Makro: Dominasi Ekstraktif vs Ekonomi Akar
Rumput}\label{konteks-makro-dominasi-ekstraktif-vs-ekonomi-akar-rumput}

Bagian ini menganalisis struktur Produk Domestik Regional Bruto (PDRB)
pada enam provinsi di Pulau Sulawesi sepanjang periode 2016--2024
menggunakan visualisasi grafik area bertumpuk (\emph{Stacked Area
Chart}). Analisis ini ditujukan untuk menguji secara empiris apakah
percepatan pertumbuhan ekonomi daerah benar-benar bersumber dari sektor
produktif masyarakat lokal atau didominasi oleh industri ekstraktif
padat modal yang mengalihkan pemanfaatan ruang dan sumber daya alam.

\begin{quote}
\textbf{Sumber Data:} Badan Pusat Statistik (BPS) Provinsi se-Sulawesi
(diolah CELIOS). Visualisasi \emph{Stacked Area Chart} memetakan
dinamika Produk Domestik Regional Bruto (PDRB) berdasarkan klasifikasi
rantai pasok hukum (\emph{Legal Supply-Chain}) untuk membandingkan
trajektori Sektor Ekstraktif, Ekonomi Akar Rumput, dan Sektor Jasa \&
Lainnya.
\end{quote}

\paragraph{A. Kerangka Dekomposisi Sektoral \& Pendekatan Rantai Pasok
Hukum (Legal
Supply-Chain)}\label{a.-kerangka-dekomposisi-sektoral-pendekatan-rantai-pasok-hukum-legal-supply-chain}

Sistem KBLI 2020 BPS membagi 17 sektor PDRB. Melalui pendekatan Legal
Supply-Chain, 17 sektor direklasifikasi menjadi 3 Klaster Makro
(Ekstraktif, Akar Rumput, Jasa). Rincian pembagian sektor, dasar
regulasi, serta intisari ketentuan hukum disajikan secara lengkap pada
\textbf{Tabel 1.1} berikut:

\subparagraph{Tabel 1.1: Reklasifikasi Sektoral PDRB KBLI 2020
Berdasarkan Pendekatan Rantai Pasok Hukum (Legal
Supply-Chain)}\label{tabel-1.1-reklasifikasi-sektoral-pdrb-kbli-2020-berdasarkan-pendekatan-rantai-pasok-hukum-legal-supply-chain}

{\def\LTcaptype{none} % do not increment counter
\begin{longtable}[]{@{}
  >{\raggedright\arraybackslash}p{(\linewidth - 8\tabcolsep) * \real{0.1905}}
  >{\raggedright\arraybackslash}p{(\linewidth - 8\tabcolsep) * \real{0.1905}}
  >{\centering\arraybackslash}p{(\linewidth - 8\tabcolsep) * \real{0.2381}}
  >{\raggedright\arraybackslash}p{(\linewidth - 8\tabcolsep) * \real{0.1905}}
  >{\raggedright\arraybackslash}p{(\linewidth - 8\tabcolsep) * \real{0.1905}}@{}}
\toprule\noalign{}
\begin{minipage}[b]{\linewidth}\raggedright
Kategori BPS
\end{minipage} & \begin{minipage}[b]{\linewidth}\raggedright
Sektor Lapangan Usaha
\end{minipage} & \begin{minipage}[b]{\linewidth}\centering
Klasifikasi Analisis
\end{minipage} & \begin{minipage}[b]{\linewidth}\raggedright
Dasar Regulasi \& Mandat Hukum
\end{minipage} & \begin{minipage}[b]{\linewidth}\raggedright
Intisari Ketentuan Hukum
\end{minipage} \\
\midrule\noalign{}
\endhead
\bottomrule\noalign{}
\endlastfoot
\textbf{Kategori B} & Pertambangan dan Penggalian & Ekstraktif & Perpres
No.~26 Tahun 2010 & Ketentuan Pasal 1 Ayat (2) mengenai pengambilan
komoditas tambang dari dalam bumi. \\
\textbf{Kategori C} & Industri Pengolahan (Smelter Logam) & Ekstraktif &
UU No.~3 Tahun 2020 \& PP No.~96 Tahun 2021 & Pasal 102--103 mewajibkan
pengolahan dan pemurnian di dalam negeri sebagai kesatuan
pertambangan. \\
\textbf{Kategori D} & Pengadaan Listrik \& Gas (PLTU Captive) &
Ekstraktif & Perpres No.~112 Tahun 2022 \& RUPTL PLN & Pasal 3 Ayat (4)
huruf b mengecualikan PLTU baru hanya bagi yang terintegrasi melayani
smelter. \\
\textbf{Kategori A} & Pertanian, Kehutanan, Perikanan & Ekonomi Akar
Rumput & KBLI 2020 BPS & Sektor pemanfaatan sumber daya hayati
terbarukan dan penyerap tenaga kerja lokal. \\
\textbf{Kategori E--U} & 13 Sektor Jasa \& Konstruksi & Sektor Jasa \&
Lainnya & Klasifikasi Standar BPS & Sektor sekunder dan tersier
penunjang perekonomian daerah. \\
\end{longtable}
}

\paragraph{B. Alur Logika Metodologis Rantai Pasok Hukum (Mengapa Kat. B
+ C + D =
Ekstraktif)}\label{b.-alur-logika-metodologis-rantai-pasok-hukum-mengapa-kat.-b-c-d-ekstraktif}

Keterkaitan ketiga kategori lapangan usaha tersebut sebagai satu
kesatuan rantai pasok ekstraktif dimodelkan dalam kerangka alur logika
hukum sebagaimana diilustrasikan pada \textbf{Bagan Alur 1.1} berikut:

\subparagraph{Bagan Alur 1.1: Alur Logika Metodologis Rantai Pasok Hukum
Sektor
Ekstraktif}\label{bagan-alur-1.1-alur-logika-metodologis-rantai-pasok-hukum-sektor-ekstraktif}

\begin{Shaded}
\begin{Highlighting}[]
\NormalTok{graph TD}
\NormalTok{    subgraph Mandat\_Smelter["2. Rantai Pasok Smelter (Kategori C)"]}
\NormalTok{        A["UU No. 3/2020 Ps. 1(1)\textless{}br/\textgreater{}\textless{}i\textgreater{}Pertambangan = Eksplorasi + Penambangan + Pengolahan/Pemurnian\textless{}/i\textgreater{}"] {-}{-}\textgreater{} B["Smelter (Industri Pengolahan / Kat. C)\textless{}br/\textgreater{}\textless{}b\textgreater{}Tahapan Wajib Pertambangan\textless{}/b\textgreater{}"]}
\NormalTok{        C["UU No. 3/2020 Ps. 102–103 \& PP 96/2021\textless{}br/\textgreater{}\textless{}i\textgreater{}Mandat Wajib Pemegang IUP Operasi Produksi\textless{}/i\textgreater{}"] {-}{-}\textgreater{} B}
\NormalTok{    end}

\NormalTok{    subgraph Mandat\_Energi["3. Rantai Pasok Energi Captive (Kategori D)"]}
\NormalTok{        D["Perpres No. 112/2022 Ps. 3(4)b\textless{}br/\textgreater{}\textless{}i\textgreater{}PLTU Baru Dilarang, KECUALI Terintegrasi Smelter\textless{}/i\textgreater{}"] {-}{-}\textgreater{} E["PLTU Captive (Pengadaan Listrik / Kat. D)\textless{}br/\textgreater{}\textless{}b\textgreater{}Instrumen Rantai Pasok Off{-}Grid\textless{}/b\textgreater{}"]}
\NormalTok{        F["RUPTL PLN 2021–2030 Hal. VI{-}24\textless{}br/\textgreater{}\textless{}i\textgreater{}Pengakuan Pasokan Khusus Smelter\textless{}/i\textgreater{}"] {-}{-}\textgreater{} E}
\NormalTok{    end}

\NormalTok{    subgraph Hulu\_Tambang["1. Sektor Hulu (Kategori B)"]}
\NormalTok{        G["Perpres No. 26/2010 Ps. 1(2)\textless{}br/\textgreater{}\textless{}b\textgreater{}Pertambangan \& Penggalian (Kat. B)\textless{}/b\textgreater{}\textless{}br/\textgreater{}\textless{}i\textgreater{}Ekstraksi SDA Tak Terbarukan\textless{}/i\textgreater{}"]}
\NormalTok{    end}

\NormalTok{    G {-}{-}\textgreater{} K["\textless{}b\textgreater{}KESIMPULAN:\textless{}/b\textgreater{}\textless{}br/\textgreater{}Kat. B + Kat. C + Kat. D = \textless{}b\textgreater{}SATU KESATUAN RANTAI PASOK EKSTRAKTIF\textless{}/b\textgreater{} yang dimandatkan hukum"]}
\NormalTok{    B {-}{-}\textgreater{} K}
\NormalTok{    E {-}{-}\textgreater{} K}
\end{Highlighting}
\end{Shaded}

\paragraph{C. Formulasi Matematis: Persamaan Agregasi Sektor Ekstraktif
(Legal Supply-Chain
Aggregation)}\label{c.-formulasi-matematis-persamaan-agregasi-sektor-ekstraktif-legal-supply-chain-aggregation}

\textbf{Persamaan Agregasi Sektor Ekstraktif (Legal Supply-Chain
Aggregation):}

\begin{Shaded}
\begin{Highlighting}[]
\NormalTok{Sektor\_Ekstraktif = PDRB(Kat.B: Pertambangan) + PDRB(Kat.C: Ind. Pengolahan) + PDRB(Kat.D: Listrik)}
\end{Highlighting}
\end{Shaded}

\emph{Keterangan Variabel:} - \texttt{Sektor\_Ekstraktif}: Total nilai
tambah bruto dari klaster industri ekstraktif yang saling terintegrasi
(Triliun Rupiah). - \texttt{PDRB(Kat.B:\ Pertambangan)}: Nilai tambah
kegiatan eksplorasi dan ekstraksi bijih mineral (BPS KBLI 2020 Kategori
B). - \texttt{PDRB(Kat.C:\ Ind.\ Pengolahan)}: Nilai tambah pemurnian
logam dasar di smelter nikel (BPS KBLI 2020 Kategori C / Golongan 24). -
\texttt{PDRB(Kat.D:\ Listrik)}: Nilai tambah penyediaan listrik batubara
khusus smelter / PLTU captive (BPS KBLI 2020 Kategori D).

\textbf{Persamaan Ekonomi Akar Rumput:}

\begin{Shaded}
\begin{Highlighting}[]
\NormalTok{Sektor\_Akar\_Rumput = PDRB(Kat.A: Pertanian, Kehutanan, dan Perikanan)}
\end{Highlighting}
\end{Shaded}

\emph{Keterangan Variabel:} - \texttt{Sektor\_Akar\_Rumput}: Nilai PDRB
pemanfaatan sumber daya hayati terbarukan (Triliun Rupiah). -
\texttt{PDRB(Kat.A)}: Agregasi nilai tambah tanaman pangan, perkebunan
rakyat, perikanan, peternakan, kehutanan.

\textbf{Persamaan Sektor Jasa \& Lainnya:}

\begin{Shaded}
\begin{Highlighting}[]
\NormalTok{Sektor\_Jasa = Jumlah PDRB (Kategori E sampai dengan Kategori U)}
\end{Highlighting}
\end{Shaded}

\emph{Keterangan Variabel:} - \texttt{Sektor\_Jasa}: Nilai tambah 13
sektor penunjang sekunder dan tersier (Triliun Rupiah). -
\texttt{PDRB(Kat.\ E\ s.d.\ U)}: Akumulasi sektor perdagangan,
konstruksi, transportasi, keuangan, pendidikan, dll.

\textbf{Persamaan Total Produk Domestik Regional Bruto (PDRB Wilayah):}

\begin{Shaded}
\begin{Highlighting}[]
\NormalTok{Total\_PDRB = Sektor\_Ekstraktif + Sektor\_Akar\_Rumput + Sektor\_Jasa}
\end{Highlighting}
\end{Shaded}

\emph{Keterangan Variabel:} - \texttt{Total\_PDRB}: Total nilai Produk
Domestik Regional Bruto wilayah atas dasar harga berlaku (Triliun
Rupiah).

\textbf{Persamaan Pangsa Kontribusi Sektor Ekstraktif (\%):}

\begin{Shaded}
\begin{Highlighting}[]
\NormalTok{Pangsa\_Ekstraktif (\%) = ( Sektor\_Ekstraktif / Total\_PDRB ) * 100}
\end{Highlighting}
\end{Shaded}

\emph{Keterangan Variabel:} - \texttt{Pangsa\_Ekstraktif\ (\%)}:
Persentase pangsa dominasi sektor ekstraktif terhadap total ekonomi
(\%).

\textbf{Persamaan Laju Pertumbuhan Tahunan Sektoral (YoY):}

\begin{Shaded}
\begin{Highlighting}[]
\NormalTok{Laju\_Pertumbuhan\_Tahunan (\%) = [ ( Nilai\_Tahun\_t {-} Nilai\_Tahun\_\{t{-}1\} ) / Nilai\_Tahun\_\{t{-}1\} ] * 100}
\end{Highlighting}
\end{Shaded}

\emph{Keterangan Variabel:} - \texttt{Laju\_Pertumbuhan\_Tahunan\ (\%)}:
Tingkat percepatan/perlambatan ekspansi tahunan sektor ekonomi (\%). -
\texttt{Nilai\_Tahun\_t}: Nilai nominal PDRB sektor pada tahun berjalan
t. - \texttt{Nilai\_Tahun\_\{t-1\}}: Nilai nominal PDRB sektor pada satu
tahun sebelumnya (t - 1).

Definisi operasional, cakupan lapangan usaha, dan institusi penyedia
data primer untuk masing-masing komponen variabel dalam sistem persamaan
di atas dipaparkan pada \textbf{Tabel 1.2} berikut:

\subparagraph{Tabel 1.2: Definisi Operasional Komponen Makroekonomi dan
Sumber Data PDRB
Sektoral}\label{tabel-1.2-definisi-operasional-komponen-makroekonomi-dan-sumber-data-pdrb-sektoral}

{\def\LTcaptype{none} % do not increment counter
\begin{longtable}[]{@{}
  >{\raggedright\arraybackslash}p{(\linewidth - 8\tabcolsep) * \real{0.1905}}
  >{\raggedright\arraybackslash}p{(\linewidth - 8\tabcolsep) * \real{0.1905}}
  >{\raggedright\arraybackslash}p{(\linewidth - 8\tabcolsep) * \real{0.1905}}
  >{\centering\arraybackslash}p{(\linewidth - 8\tabcolsep) * \real{0.2381}}
  >{\raggedright\arraybackslash}p{(\linewidth - 8\tabcolsep) * \real{0.1905}}@{}}
\toprule\noalign{}
\begin{minipage}[b]{\linewidth}\raggedright
Komponen Analisis
\end{minipage} & \begin{minipage}[b]{\linewidth}\raggedright
Cakupan Lapangan Usaha
\end{minipage} & \begin{minipage}[b]{\linewidth}\raggedright
Definisi Operasional
\end{minipage} & \begin{minipage}[b]{\linewidth}\centering
Satuan Nilai
\end{minipage} & \begin{minipage}[b]{\linewidth}\raggedright
Sumber Data Primer
\end{minipage} \\
\midrule\noalign{}
\endhead
\bottomrule\noalign{}
\endlastfoot
\textbf{Sektor Ekstraktif} & Kategori B, Kategori C, Kategori D &
Akumulasi nilai tambah pertambangan nikel, smelter logam dasar, dan PLTU
captive. & Triliun Rupiah & BPS Provinsi (SIMDASI) \\
\textbf{Ekonomi Akar Rumput} & Kategori A & Nilai tambah pertanian,
perkebunan, kehutanan, dan perikanan. & Triliun Rupiah & BPS Provinsi \\
\textbf{Sektor Jasa \& Lainnya} & Kategori E hingga U & Nilai tambah
gabungan perdagangan, konstruksi, transportasi, keuangan, dan jasa. &
Triliun Rupiah & BPS Provinsi \\
\textbf{Total PDRB Wilayah} & Seluruh 17 Kategori & Total nilai PDRB
wilayah atas dasar harga berlaku pada tahun berjalan. & Triliun Rupiah &
BPS Provinsi \\
\textbf{Pangsa Ekstraktif (\%)} & Rasio Kontribusi & Persentase
kontribusi sektor ekstraktif terhadap total perekonomian. & Persen (\%)
& Hasil Olahan CELIOS \\
\end{longtable}
}

\paragraph{D. Analisis Temuan Empiris: Ketimpangan Struktural Sulawesi
Tengah}\label{d.-analisis-temuan-empiris-ketimpangan-struktural-sulawesi-tengah}

Penerapan formulasi di atas menunjukkan bahwa di \textbf{Sulawesi Tengah
(sebagai pusat hilirisasi)}, ekspansi industri ekstraktif menguasai
\textbf{55.8\% dari total PDRB provinsi} pada tahun 2024, memperlihatkan
dominasi yang sangat kuat dibanding provinsi lainnya.

\subsubsection{1.1.2 Pemusatan Sektor Ekstraktif di Kabupaten
se-Sulawesi
Tengah}\label{pemusatan-sektor-ekstraktif-di-kabupaten-se-sulawesi-tengah}

Jika dianalisis secara spasial pada tingkat kabupaten di Sulawesi
Tengah, terlihat konsentrasi kegiatan industri ekstraktif. Kabupaten
\textbf{Morowali} dan \textbf{Morowali Utara} mendominasi struktur PDRB
provinsi melalui pengembangan kawasan industri hilirisasi dan PLTU
Captive. Analisis ini membandingkan komposisi ketiga sektor advokatif di
seluruh 13 kabupaten/kota se-Sulawesi Tengah pada tahun terbaru (2024).

\begin{quote}
\textbf{Sumber Data:} Badan Pusat Statistik (BPS) Kabupaten se-Sulawesi
Tengah (diolah CELIOS). Visualisasi \emph{Stacked Bar Chart} memetakan
struktur Produk Domestik Regional Bruto (PDRB) tahun 2024 pada seluruh
13 kabupaten/kota untuk mengidentifikasi tingkat konsentrasi sektoral
dan polarisasi spasial antara sentra industri pengolahan nikel dengan
daerah non-sentra.
\end{quote}

\paragraph{A. Rasionalitas Spasial \& Urgensi Dekomposisi Sektoral
Tingkat
Kabupaten}\label{a.-rasionalitas-spasial-urgensi-dekomposisi-sektoral-tingkat-kabupaten}

Analisis agregat pada tingkat provinsi sering kali menghasilkan
\textbf{Bias Ilusi Agregat (Aggregate Illusion Bias)}, di mana angka
pertumbuhan ekonomi makro yang tinggi memberi kesan seolah seluruh
wilayah menikmati kemakmuran yang seimbang. Namun, ketika data
didekomposisi ke tingkat kabupaten/kota, terlihat jurang pemisah ekonomi
yang sangat tajam antara wilayah \textbf{Enklave Industri Ekstraktif}
dengan daerah agraris tradisional sekitarnya.

\paragraph{B. Alur Logika Analisis Disparitas Spasial
Kabupaten}\label{b.-alur-logika-analisis-disparitas-spasial-kabupaten}

Kerangka kerja metodologis dalam membedah ketimpangan intra-provinsial
ini diilustrasikan pada \textbf{Bagan Alur 1.2} berikut:

\subparagraph{Bagan Alur 1.2: Alur Logika Metodologis Dekomposisi
Spasial PDRB Tingkat Kabupaten se-Sulawesi
Tengah}\label{bagan-alur-1.2-alur-logika-metodologis-dekomposisi-spasial-pdrb-tingkat-kabupaten-se-sulawesi-tengah}

\begin{Shaded}
\begin{Highlighting}[]
\NormalTok{graph TD}
\NormalTok{    subgraph Data\_BPS["1. Input Data Statistik Daerah"]}
\NormalTok{        A["BPS Kabupaten se{-}Sulawesi Tengah\textless{}br/\textgreater{}\textless{}i\textgreater{}PDRB 17 Sektor Lapangan Usaha ADHB\textless{}/i\textgreater{}"]}
\NormalTok{    end}

\NormalTok{    subgraph Reklasifikasi["2. Reklasifikasi Legal Supply{-}Chain"]}
\NormalTok{        B["Ekstraktif = Kat. B + Kat. C + Kat. D"]}
\NormalTok{        C["Akar Rumput = Kat. A"]}
\NormalTok{        D["Jasa \& Lainnya = Kat. E s.d. U"]}
\NormalTok{    end}

\NormalTok{    subgraph Analisis\_Disparitas["3. Output Evaluasi Spasial"]}
\NormalTok{        E["\textless{}b\textgreater{}Sentra Hilirisasi (Enclave Industri):\textless{}/b\textgreater{}\textless{}br/\textgreater{}Morowali \& Morut menguasai Sektor Ekstraktif Tertinggi"]}
\NormalTok{        F["\textless{}b\textgreater{}Non{-}Sentra (Pertanian Rakyat):\textless{}/b\textgreater{}\textless{}br/\textgreater{}11 Kabupaten tertinggal (\textless{}11\% Porsi Ekstraktif)"]}
\NormalTok{    end}

\NormalTok{    A {-}{-}\textgreater{} B}
\NormalTok{    A {-}{-}\textgreater{} C}
\NormalTok{    A {-}{-}\textgreater{} D}
\NormalTok{    B {-}{-}\textgreater{} E}
\NormalTok{    C {-}{-}\textgreater{} F}
\NormalTok{    D {-}{-}\textgreater{} F}
\end{Highlighting}
\end{Shaded}

\paragraph{C. Formulasi Matematis: Persamaan Agregasi Sektoral Kabupaten
(Legal Supply-Chain
Aggregation)}\label{c.-formulasi-matematis-persamaan-agregasi-sektoral-kabupaten-legal-supply-chain-aggregation}

\textbf{Persamaan Agregasi Sektor Ekstraktif Tingkat Kabupaten:}

\begin{Shaded}
\begin{Highlighting}[]
\NormalTok{Sektor\_Ekstraktif\_Kabupaten = PDRB\_Kab(Kat.B: Pertambangan) + PDRB\_Kab(Kat.C: Ind. Pengolahan) + PDRB\_Kab(Kat.D: Listrik)}
\end{Highlighting}
\end{Shaded}

\emph{Keterangan Variabel:} - \texttt{Sektor\_Ekstraktif\_Kabupaten}:
Total nilai tambah sektor ekstraktif di tingkat kabupaten target
(satuan: Triliun Rupiah). - \texttt{PDRB\_Kab(Kat.B:\ Pertambangan)}:
Nilai PDRB kabupaten dari aktivitas penambangan bijih logam dan galian
(BPS Kategori B). - \texttt{PDRB\_Kab(Kat.C:\ Ind.\ Pengolahan)}: Nilai
PDRB kabupaten dari industri peleburan logam dasar / smelter (BPS
Kategori C). - \texttt{PDRB\_Kab(Kat.D:\ Listrik)}: Nilai PDRB kabupaten
dari penyediaan daya listrik batubara captive (BPS Kategori D).

\textbf{Persamaan Total Produk Domestik Regional Bruto Tingkat
Kabupaten:}

\begin{Shaded}
\begin{Highlighting}[]
\NormalTok{Total\_PDRB\_Kabupaten = Sektor\_Ekstraktif\_Kabupaten + Sektor\_Akar\_Rumput\_Kabupaten + Sektor\_Jasa\_Kabupaten}
\end{Highlighting}
\end{Shaded}

\emph{Keterangan Variabel:} - \texttt{Total\_PDRB\_Kabupaten}: Total
output perekonomian bruto kabupaten target atas dasar harga berlaku
(satuan: Triliun Rupiah). - \texttt{Sektor\_Ekstraktif\_Kabupaten}:
Nilai tambah bruto sektor ekstraktif terintegrasi di kabupaten (Triliun
Rupiah). - \texttt{Sektor\_Akar\_Rumput\_Kabupaten}: Nilai tambah sektor
pertanian, kehutanan, dan perikanan di kabupaten (Triliun Rupiah). -
\texttt{Sektor\_Jasa\_Kabupaten}: Nilai tambah sektor perdagangan,
transportasi, dan jasa layanan di kabupaten (Triliun Rupiah).

\textbf{Persamaan Porsi Sektoral dalam Kabupaten (Porsi (\%) pada
Tooltip Dashboard):}

\begin{Shaded}
\begin{Highlighting}[]
\NormalTok{Porsi\_Sektor\_Kabupaten (\%) = ( Nilai\_Sektor\_Kabupaten / Total\_PDRB\_Kabupaten ) * 100}
\end{Highlighting}
\end{Shaded}

\emph{Keterangan Variabel:} - \texttt{Porsi\_Ekstraktif\ (\%)}:
Persentase kontribusi Sektor Ekstraktif: ( Sektor\_Ekstraktif /
Total\_PDRB ) * 100 (misal Morowali: 45.2\%). -
\texttt{Porsi\_Jasa\ (\%)}: Persentase kontribusi Sektor Jasa \&
Lainnya: ( Sektor\_Jasa / Total\_PDRB ) * 100 (misal Morowali: 54.0\%).
- \texttt{Porsi\_Akar\_Rumput\ (\%)}: Persentase kontribusi Sektor
Ekonomi Akar Rumput: ( Sektor\_Akar\_Rumput / Total\_PDRB ) * 100 (misal
Morowali: 0.8\%). - \texttt{Total\_PDRB\_Kabupaten}: Total nilai nominal
PDRB seluruh sektor di kabupaten target (Triliun Rupiah).

\paragraph{D. Rincian Definisi Operasional \& Matriks Distribusi PDRB 13
Kabupaten}\label{d.-rincian-definisi-operasional-matriks-distribusi-pdrb-13-kabupaten}

Penerapan sistem persamaan di atas terhadap seluruh 13 kabupaten dan
kota di Provinsi Sulawesi Tengah pada tahun 2024 disajikan secara
komprehensif pada \textbf{Tabel 1.3} berikut:

\subparagraph{Tabel 1.3: Distribusi Nilai Tambah Bruto dan Komposisi
Sektoral PDRB 13 Kabupaten/Kota di Sulawesi Tengah (Tahun
2024)}\label{tabel-1.3-distribusi-nilai-tambah-bruto-dan-komposisi-sektoral-pdrb-13-kabupatenkota-di-sulawesi-tengah-tahun-2024}

{\def\LTcaptype{none} % do not increment counter
\begin{longtable}[]{@{}
  >{\raggedright\arraybackslash}p{(\linewidth - 16\tabcolsep) * \real{0.0930}}
  >{\centering\arraybackslash}p{(\linewidth - 16\tabcolsep) * \real{0.1163}}
  >{\centering\arraybackslash}p{(\linewidth - 16\tabcolsep) * \real{0.1163}}
  >{\centering\arraybackslash}p{(\linewidth - 16\tabcolsep) * \real{0.1163}}
  >{\centering\arraybackslash}p{(\linewidth - 16\tabcolsep) * \real{0.1163}}
  >{\centering\arraybackslash}p{(\linewidth - 16\tabcolsep) * \real{0.1163}}
  >{\centering\arraybackslash}p{(\linewidth - 16\tabcolsep) * \real{0.1163}}
  >{\centering\arraybackslash}p{(\linewidth - 16\tabcolsep) * \real{0.1163}}
  >{\raggedright\arraybackslash}p{(\linewidth - 16\tabcolsep) * \real{0.0930}}@{}}
\toprule\noalign{}
\begin{minipage}[b]{\linewidth}\raggedright
Kabupaten / Kota
\end{minipage} & \begin{minipage}[b]{\linewidth}\centering
Akar Rumput (T Rp)
\end{minipage} & \begin{minipage}[b]{\linewidth}\centering
Ekstraktif (T Rp)
\end{minipage} & \begin{minipage}[b]{\linewidth}\centering
Jasa (T Rp)
\end{minipage} & \begin{minipage}[b]{\linewidth}\centering
Total PDRB (T Rp)
\end{minipage} & \begin{minipage}[b]{\linewidth}\centering
Porsi Akar Rumput (\%)
\end{minipage} & \begin{minipage}[b]{\linewidth}\centering
Porsi Ekstraktif (\%)
\end{minipage} & \begin{minipage}[b]{\linewidth}\centering
Porsi Jasa (\%)
\end{minipage} & \begin{minipage}[b]{\linewidth}\raggedright
Basis Utama Ekonomi
\end{minipage} \\
\midrule\noalign{}
\endhead
\bottomrule\noalign{}
\endlastfoot
\textbf{Morowali} & 2.70 & 157.17 & 187.85 & \textbf{347.72} & 0.8\% &
45.2\% & 54.0\% & Hilirisasi Nikel (Smelter \& PLTU) \\
\textbf{Banggai} & 8.85 & 20.63 & 51.99 & \textbf{81.47} & 10.9\% &
25.3\% & 63.8\% & Migas, Tambang \& Perdagangan \\
\textbf{Palu} & 1.24 & 4.56 & 60.03 & \textbf{65.84} & 1.9\% & 6.9\% &
91.2\% & Jasa, Perdagangan \& Pemerintahan \\
\textbf{Morowali Utara} & 5.17 & 19.22 & 36.08 & \textbf{60.47} & 8.5\%
& 31.8\% & 59.7\% & Hilirisasi Nikel (Smelter GNI) \\
\textbf{Parigi Moutong} & 9.97 & 1.93 & 35.05 & \textbf{46.95} & 21.2\%
& 4.1\% & 74.7\% & Pertanian Pangan \& Hortikultura \\
\textbf{Donggala} & 5.96 & 3.53 & 23.57 & \textbf{33.05} & 18.0\% &
10.7\% & 71.3\% & Pertanian, Perkebunan \& Galian C \\
\textbf{Poso} & 4.96 & 0.40 & 20.12 & \textbf{25.48} & 19.5\% & 1.6\% &
79.0\% & Pertanian \& Perkebunan Kakao \\
\textbf{Sigi} & 5.17 & 0.83 & 19.13 & \textbf{25.12} & 20.6\% & 3.3\% &
76.1\% & Pertanian Pangan \& Hortikultura \\
\textbf{Toli-Toli} & 4.35 & 0.44 & 17.36 & \textbf{22.15} & 19.7\% &
2.0\% & 78.4\% & Perkebunan Cengkeh \& Perikanan \\
\textbf{Buol} & 3.77 & 1.15 & 10.67 & \textbf{15.58} & 24.2\% & 7.4\% &
68.5\% & Kelapa Sawit \& Tanaman Pangan \\
\textbf{Tojo Una-Una} & 2.88 & 0.61 & 11.21 & \textbf{14.71} & 19.6\% &
4.2\% & 76.2\% & Pertanian \& Pariwisata Bahari \\
\textbf{Banggai Kepulauan} & 2.53 & 0.18 & 8.04 & \textbf{10.75} &
23.5\% & 1.7\% & 74.8\% & Perikanan Tangkap \& Kelautan \\
\textbf{Banggai Laut} & 1.80 & 0.12 & 4.52 & \textbf{6.45} & 27.9\% &
1.9\% & 70.2\% & Perikanan \& Budidaya Laut \\
\end{longtable}
}

\paragraph{E. Analisis Temuan Empiris: Polarisasi Ekstrem Morowali vs
Daerah
Non-Smelter}\label{e.-analisis-temuan-empiris-polarisasi-ekstrem-morowali-vs-daerah-non-smelter}

Data empiris pada Tabel 1.3 mengungkap bukti polarisasi ekonomi wilayah
yang sangat ekstrem di Sulawesi Tengah:

\begin{enumerate}
\def\labelenumi{\arabic{enumi}.}
\tightlist
\item
  \textbf{Dominasi Sektor Ekstraktif Morowali:} Kabupaten Morowali
  mencatatkan nilai sektor ekstraktif sebesar Rp 157.17 Triliun atau
  menguasai porsi 45.2\% dari total kue ekonomi kabupatennya (Rp 347.72
  Triliun). Nilai sektor ekstraktif Morowali saja melampaui gabungan
  total PDRB dari delapan kabupaten lainnya di Sulawesi Tengah.
\item
  \textbf{Pemusatan pada Dua Sentra Hilirisasi:} Kabupaten Morowali dan
  Morowali Utara merupakan dua daerah dengan nilai Sektor Ekstraktif
  tertinggi di Sulawesi Tengah, membuktikan bahwa percepatan output
  industri pertambangan dan hilirisasi terkunci pada kawasan industri
  smelter.
\item
  \textbf{Ketertinggalan Daerah Non-Sentra:} Sebaliknya, delapan
  kabupaten lainnya (seperti Banggai Laut, Banggai Kepulauan, Tojo
  Una-Una, Buol, Toli-Toli, Sigi, Poso, dan Donggala) memiliki porsi
  Sektor Ekstraktif yang sangat rendah (\textless11\%) dan tetap
  bergantung pada sektor pertanian rakyat (Akar Rumput) berproduktivitas
  rendah dengan keterbatasan akses terhadap nilai tambah modal.
\end{enumerate}

\subsubsection{1.1.3 Perbandingan Distribusi 17 Sektor Komoditas per
Provinsi (Small Multiples, Tahun
Terbaru)}\label{perbandingan-distribusi-17-sektor-komoditas-per-provinsi-small-multiples-tahun-terbaru}

Visualisasi komparatif \textbf{Small Multiples Horizontal Bar Chart}
membedah struktur 17 sektor lapangan usaha KBLI 2020 secara terpisah
pada enam provinsi di Pulau Sulawesi pada tahun terbaru (2024). Setiap
panel provinsi menampilkan sektor yang diurutkan dari penyumbang
terbesar hingga terkecil dengan skala sumbu nilai yang disetarakan
secara seragam untuk memastikan validitas komparasi lintas wilayah.

\begin{quote}
\textbf{Sumber Data Resmi \& Deskripsi Visualisasi:} Badan Pusat
Statistik (BPS) Provinsi se-Sulawesi (diolah CELIOS). Visualisasi
\emph{Small Multiples Horizontal Bar Chart} menyajikan dekomposisi 17
sektor PDRB tahun 2024 di 6 provinsi se-Pulau Sulawesi. Sumbu X
disetarakan pada rentang nilai seragam ({[}0, 178.5 Triliun Rp{]})
dengan pewarnaan berdasarkan 3 klaster makro (Merah: Ekstraktif, Hijau:
Ekonomi Akar Rumput, Abu-abu: Sektor Jasa \& Lainnya) guna
mengidentifikasi spesialisasi dan anomali struktural ekonomi
masing-masing provinsi.
\end{quote}

\paragraph{A. Kerangka Konseptual \& Standardisasi Skala Komparatif
(Uniform Scale Small
Multiples)}\label{a.-kerangka-konseptual-standardisasi-skala-komparatif-uniform-scale-small-multiples}

Dalam analisis data multidimensi lintas wilayah, penggunaan skala
dinamis mandiri (\emph{independent dynamic scaling}) pada masing-masing
panel sering kali menimbulkan \textbf{Bias Distorsi Visual Komparatif
(Visual Comparison Bias)}. Tanpa penyetaraan batas skala maksimum,
sektor dengan nominal kecil di provinsi ber-PDRB rendah dapat terlihat
secara visual setara dengan sektor bernilai ratusan triliun di provinsi
ber-PDRB besar. Oleh karena itu, metodologi ini menetapkan batas skala
maksimum sumbu X yang seragam (\emph{Uniform Scale Bound}) sebesar nilai
maksimum sektor tertinggi di seluruh pulau ditambah faktor ruang margin
sebesar 15\%.

\paragraph{B. Alur Logika Metodologis Analisis Small Multiples 17
Sektor}\label{b.-alur-logika-metodologis-analisis-small-multiples-17-sektor}

Kerangka operasionalisasi analisis perbandingan terpisah 17 sektor
lapangan usaha ini dimodelkan dalam kerangka alur logika sebagaimana
diilustrasikan pada \textbf{Bagan Alur 1.3} berikut:

\subparagraph{Bagan Alur 1.3: Alur Logika Metodologis Analisis
Komparatif Small Multiples 17 Sektor PDRB per
Provinsi}\label{bagan-alur-1.3-alur-logika-metodologis-analisis-komparatif-small-multiples-17-sektor-pdrb-per-provinsi}

\begin{Shaded}
\begin{Highlighting}[]
\NormalTok{graph TD}
\NormalTok{    subgraph Input\_Data["1. Input Data Statistik Provinsi"]}
\NormalTok{        A["BPS Provinsi se{-}Sulawesi\textless{}br/\textgreater{}\textless{}i\textgreater{}PDRB 17 Sektor Lapangan Usaha ADHB\textless{}/i\textgreater{}"]}
\NormalTok{    end}

\NormalTok{    subgraph Standardisasi["2. Pemrosesan \& Standardisasi Data"]}
\NormalTok{        B["Konversi Nilai: Triliun Rp = Miliar Rp / 1000"]}
\NormalTok{        C["Porsi Sektor (\%) = (Nilai Sektor / Total PDRB) * 100"]}
\NormalTok{        D["Skala X Seragam: [0, max(X) * 1.15]\textless{}br/\textgreater{}\textless{}i\textgreater{}Mencegah Visual Comparison Bias\textless{}/i\textgreater{}"]}
\NormalTok{        E["Pewarnaan 3 Klaster Kritis:\textless{}br/\textgreater{}• Merah (Kat. B,C,D: Ekstraktif)\textless{}br/\textgreater{}• Hijau (Kat. A: Akar Rumput)\textless{}br/\textgreater{}• Abu{-}abu (Kat. E–U: Jasa \& Lainnya)"]}
\NormalTok{    end}

\NormalTok{    subgraph Komparasi\_Regional["3. Output Komparasi Antarwilayah"]}
\NormalTok{        F["\textless{}b\textgreater{}Sulawesi Tengah \& Sultra:\textless{}/b\textgreater{}\textless{}br/\textgreater{}Anomali Lonjakan Sektor Ekstraktif (Smelter \& Tambang)"]}
\NormalTok{        G["\textless{}b\textgreater{}Sulsel, Sulbar, Gorontalo \& Sulut:\textless{}/b\textgreater{}\textless{}br/\textgreater{}Struktur Perekonomian Bertumpu pada Pertanian \& Jasa"]}
\NormalTok{    end}

\NormalTok{    A {-}{-}\textgreater{} B}
\NormalTok{    B {-}{-}\textgreater{} C}
\NormalTok{    B {-}{-}\textgreater{} D}
\NormalTok{    C {-}{-}\textgreater{} E}
\NormalTok{    D {-}{-}\textgreater{} E}
\NormalTok{    E {-}{-}\textgreater{} F}
\NormalTok{    E {-}{-}\textgreater{} G}
\end{Highlighting}
\end{Shaded}

\paragraph{C. Formulasi Matematis: Persamaan Agregasi dan Porsi 17
Sektor
Komoditas}\label{c.-formulasi-matematis-persamaan-agregasi-dan-porsi-17-sektor-komoditas}

Kalkulasi perbandingan sektoral dan normalisasi skala grafik dihitung
menggunakan sistem formulasi berikut:

\textbf{Persamaan Normalisasi Nilai Sektor ke Satuan Triliun Rupiah:}

\begin{Shaded}
\begin{Highlighting}[]
\NormalTok{Nilai\_Sektor\_Triliun = Nilai\_Sektor\_Miliar / 1000}
\end{Highlighting}
\end{Shaded}

\emph{Keterangan Variabel:} - \texttt{Nilai\_Sektor\_Triliun}: Nilai
tambah bruto sektor lapangan usaha dalam satuan baku Triliun Rupiah. -
\texttt{Nilai\_Sektor\_Miliar}: Nilai nominal PDRB mentah dari publikasi
resmi BPS (satuan: Miliar Rupiah).

\textbf{Persamaan Porsi Sektoral per Provinsi (Porsi (\%) pada Tooltip
Dashboard):}

\begin{Shaded}
\begin{Highlighting}[]
\NormalTok{Porsi\_Sektor\_Provinsi (\%) = ( Nilai\_Sektor\_Provinsi / Total\_PDRB\_Provinsi ) * 100}
\end{Highlighting}
\end{Shaded}

\emph{Keterangan Variabel:} - \texttt{Porsi\_Sektor\_Provinsi\ (\%)}:
Persentase kontribusi sektor target terhadap keseluruhan total PDRB
provinsi bersangkutan (satuan: Persen / \%). Angka ini ditampilkan pada
tooltip `Porsi (\%)' di dashboard. - \texttt{Nilai\_Sektor\_Provinsi}:
Nilai tambah bruto sektor lapangan usaha target di provinsi bersangkutan
(Triliun Rupiah). - \texttt{Total\_PDRB\_Provinsi}: Total nilai nominal
PDRB seluruh 17 sektor di provinsi bersangkutan (Triliun Rupiah).

\textbf{Persamaan Batas Maksimum Skala Sumbu X Seragam (Uniform Scale
Bound):}

\begin{Shaded}
\begin{Highlighting}[]
\NormalTok{X\_Max\_Seragam = Maksimum(Seluruh\_Nilai\_Sektor\_Semua\_Provinsi) * 1.15}
\end{Highlighting}
\end{Shaded}

\emph{Keterangan Variabel:} - \texttt{X\_Max\_Seragam}: Nilai batas atas
sumbu X yang diaplikasikan secara identik pada semua grafik Small
Multiples. - \texttt{Maksimum(...)}: Sektor dengan nominal PDRB terbesar
di seluruh pulau (contoh: Sektor Pertanian di Sulsel). - \texttt{1.15}:
Faktor pengali batas margin (+15\%) untuk ruang keterangan (label
space).

\paragraph{D. Rincian Data Empiris: Matriks Perbandingan Sektor Unggulan
6
Provinsi}\label{d.-rincian-data-empiris-matriks-perbandingan-sektor-unggulan-6-provinsi}

Penerapan sistem perbandingan komparatif di atas mengidentifikasi 5
sektor tulang punggung utama (\emph{top 5 contributors}) dari 17
Kategori BPS di tiap provinsi pada tahun 2024, sebagaimana dirinci pada
\textbf{Tabel 1.4} berikut:

\subparagraph{Tabel 1.4: 5 Sektor Lapangan Usaha Penyumbang Utama PDRB
di 6 Provinsi Sulawesi (Tahun
2024)}\label{tabel-1.4-5-sektor-lapangan-usaha-penyumbang-utama-pdrb-di-6-provinsi-sulawesi-tahun-2024}

{\def\LTcaptype{none} % do not increment counter
\begin{longtable}[]{@{}
  >{\raggedright\arraybackslash}p{(\linewidth - 10\tabcolsep) * \real{0.1667}}
  >{\raggedright\arraybackslash}p{(\linewidth - 10\tabcolsep) * \real{0.1667}}
  >{\raggedright\arraybackslash}p{(\linewidth - 10\tabcolsep) * \real{0.1667}}
  >{\raggedright\arraybackslash}p{(\linewidth - 10\tabcolsep) * \real{0.1667}}
  >{\raggedright\arraybackslash}p{(\linewidth - 10\tabcolsep) * \real{0.1667}}
  >{\raggedright\arraybackslash}p{(\linewidth - 10\tabcolsep) * \real{0.1667}}@{}}
\toprule\noalign{}
\begin{minipage}[b]{\linewidth}\raggedright
Provinsi
\end{minipage} & \begin{minipage}[b]{\linewidth}\raggedright
Peringkat 1
\end{minipage} & \begin{minipage}[b]{\linewidth}\raggedright
Peringkat 2
\end{minipage} & \begin{minipage}[b]{\linewidth}\raggedright
Peringkat 3
\end{minipage} & \begin{minipage}[b]{\linewidth}\raggedright
Peringkat 4
\end{minipage} & \begin{minipage}[b]{\linewidth}\raggedright
Peringkat 5
\end{minipage} \\
\midrule\noalign{}
\endhead
\bottomrule\noalign{}
\endlastfoot
\textbf{Sulawesi Tengah} & \textbf{376.95} & Industri Pengolahan &
41.2\% & Pertanian \& Perikanan & 15.8\% \\
\textbf{Sulawesi Tenggara} & \textbf{189.48} & Pertanian \& Perikanan &
23.5\% & Pertambangan \& Penggalian & 21.1\% \\
\textbf{Sulawesi Selatan} & \textbf{696.25} & Pertanian \& Perikanan &
21.8\% & Perdagangan Besar \& Eceran & 14.8\% \\
\textbf{Sulawesi Utara} & \textbf{187.37} & Pertanian \& Perikanan &
20.6\% & Perdagangan Besar \& Eceran & 13.6\% \\
\textbf{Sulawesi Barat} & \textbf{64.21} & Pertanian \& Perikanan &
46.1\% & Industri Pengolahan & 10.4\% \\
\textbf{Gorontalo} & \textbf{54.56} & Pertanian \& Perikanan & 37.3\% &
Perdagangan Besar \& Eceran & 14.2\% \\
\end{longtable}
}

\paragraph{E. Analisis Temuan Empiris \& Interpretasi Sektoral
Dashboard}\label{e.-analisis-temuan-empiris-interpretasi-sektoral-dashboard}

Visualisasi \emph{Small Multiples} dan matriks Tabel 1.4 memvalidasi
hipotesis adanya \textbf{Disparitas Struktural} yang dipicu oleh
ekspansi industri nikel:

\begin{enumerate}
\def\labelenumi{\arabic{enumi}.}
\tightlist
\item
  \textbf{Dominasi Absolut Ekstraktif di Sentra Smelter:} Sulawesi
  Tengah dan Sulawesi Tenggara menunjukkan pola yang identik, di mana
  Sektor Ekstraktif (Industri Pengolahan Logam Dasar dan Pertambangan)
  menjadi tulang punggung absolut. Khususnya di Sulawesi Tengah, jarak
  nilai (gap) antara sektor ekstraktif dengan sektor lainnya sangat
  ekstrem.
\item
  \textbf{Perekonomian Berbasis Akar Rumput di Provinsi Lain:} Empat
  provinsi lainnya (Sulawesi Selatan, Sulawesi Barat, Gorontalo, dan
  Sulawesi Utara) tetap mengandalkan Sektor Pertanian (Kategori A)
  sebagai penyumbang terbesar PDRB, mencerminkan resiliensi ekonomi akar
  rumput di wilayah non-smelter.
\item
  \textbf{Urgensi Normalisasi Skala Visual:} Penggunaan skala X seragam
  ({[}0, 168 Triliun Rp{]}) memastikan pembaca menyadari bahwa meskipun
  Sektor Industri Pengolahan menempati Peringkat 1 di Sulawesi Tengah,
  nominal absolutnya belum tentu setara dengan Sektor Pertanian di
  Sulawesi Selatan. Visualisasi terstandarisasi ini mencegah kesimpulan
  spekulatif.
\end{enumerate}

\begin{center}\rule{0.5\linewidth}{0.5pt}\end{center}

\subsection{1.2 Konsentrasi Kawasan Industri \& PLTU
Captive}\label{konsentrasi-kawasan-industri-pltu-captive}

Intensifikasi industri pengolahan mineral di Pulau Sulawesi berpusat
pada pembangunan mega-smelter yang ditopang secara mutlak oleh
pembangkit listrik tenaga uap khusus (\emph{PLTU Captive}) batu bara
non-jaringan (\emph{off-grid}). Bagian ini mengombinasikan
\textbf{Analisis Spasial Deskriptif} untuk mengidentifikasi tingkat
pemusatan fasilitas dan kapasitas energi fosil di enam provinsi, dengan
\textbf{Uji Tabulasi Silang Panel (Inferential Spatiotemporal
Crosstabulation)} berstandar SPSS guna membuktikan secara ilmiah
keterkaitan antara ekspansi PLTU Captive terhadap kehilangan tutupan
hutan di Pulau Sulawesi.

\begin{quote}
\textbf{Sumber Data Resmi \& Deskripsi Metodologis:} Kementerian Energi
dan Sumber Daya Mineral (ESDM / Minerbaone), Global Energy Monitor (GEM
Coal Plant Tracker), dan Global Forest Watch (GFW / University of
Maryland) (diolah CELIOS). Visualisasi \emph{Bar Chart} Konsentrasi
Industri dan Pemetaan Spasial menyajikan distribusi 778 unit fasilitas
smelter serta 9,825 MW kapasitas terpasang aktif PLTU captive di 6
provinsi se-Pulau Sulawesi. Analisis dipadukan dengan Uji Tabulasi
Silang Data Panel Spasiotemporal (Chi-Square Test \& Risk Odds Ratio,
N=60) untuk menguji keterkaitan ekspansi energi fosil industri terhadap
eskalasi deforestasi komoditas.
\end{quote}

\subsubsection{A. Pemusatan Spasial Fasilitas Smelter dan PLTU
Captive}\label{a.-pemusatan-spasial-fasilitas-smelter-dan-pltu-captive}

Intensifikasi industri pengolahan nikel di Sulawesi berpusat pada
fasilitas mega-smelter. Pengoperasian \textbf{778 fasilitas smelter}
didukung oleh kapasitas energi batu bara \textbf{9,825 MW dari PLTU
Captive}. Berbeda dengan sistem kelistrikan umum PLN, pembangkit ini
dikembangkan secara internal untuk menyokong operasi kawasan industri.

\subsubsection{B. Metodologi: Analisis Spasial \& Uji Tabulasi
Silang}\label{b.-metodologi-analisis-spasial-uji-tabulasi-silang}

Pengujian keterkaitan antara pembangunan PLTU Captive dengan kehilangan
tutupan hutan dioperasionalkan melalui Standar Operasional Prosedur
(SOP) tabulasi silang berstandar SPSS. Rangkaian tahapan logika
metodologis, asumsi frekuensi harapan, hingga estimasi faktor risiko
dimodelkan pada \textbf{Bagan Alur 1.4} berikut:

\subparagraph{Bagan Alur 1.4: Standar Operasional Prosedur (SOP) \& Alur
Logika Uji Tabulasi Silang (Crosstab) PLTU Captive vs
Deforestasi}\label{bagan-alur-1.4-standar-operasional-prosedur-sop-alur-logika-uji-tabulasi-silang-crosstab-pltu-captive-vs-deforestasi}

\begin{Shaded}
\begin{Highlighting}[]
\NormalTok{flowchart TD}
\NormalTok{    A(["Start: Input Data"]) {-}{-}\textgreater{} B\{"Apakah Data\textless{}br/\textgreater{}Kategorikal?"\}}
    
\NormalTok{    B {-}{-} TIDAK (Numerik) {-}{-}\textgreater{} C["Lakukan Diskritisasi/Binning\textless{}br/\textgreater{}Ubah Angka jadi Kategori\textless{}br/\textgreater{}Misal: Tinggi vs Rendah"]}
\NormalTok{    C {-}{-}\textgreater{} D}
    
\NormalTok{    B {-}{-} YA {-}{-}\textgreater{} D["Penanganan Missing Values\textless{}br/\textgreater{}SPSS: Listwise Deletion\textless{}br/\textgreater{}Hapus baris yang datanya bolong"]}
    
\NormalTok{    D {-}{-}\textgreater{} E["Jalankan Perhitungan\textless{}br/\textgreater{}Tabel Crosstab 2x2"]}
    
\NormalTok{    E {-}{-}\textgreater{} F\{"Cek Asumsi SPSS:\textless{}br/\textgreater{}Expected Count \textgreater{}= 5 ?"\}}
    
\NormalTok{    F {-}{-} TIDAK (Ada sel bernilai 0 \textless{}br/\textgreater{}atau \textgreater{}20\% sel nilainya \textless{}5) {-}{-}\textgreater{} G["Asumsi Pearson Chi{-}Square\textless{}br/\textgreater{}DILANGGAR / CACAT"]}
\NormalTok{    G {-}{-}\textgreater{} H["Gunakan Uji Alternatif:\textless{}br/\textgreater{}Fisher\textquotesingle{}s Exact Test"]}
\NormalTok{    H {-}{-}\textgreater{} K}
    
\NormalTok{    F {-}{-} YA (Memenuhi Syarat) {-}{-}\textgreater{} I["Asumsi Terpenuhi\textless{}br/\textgreater{}Baca Nilai Pearson Chi{-}Square"]}
\NormalTok{    I {-}{-}\textgreater{} J\{"Cek P{-}Value\textless{}br/\textgreater{}(Asymp. Sig)"\}}
    
\NormalTok{    J {-}{-} P{-}Value \textgreater{}= 0.05 {-}{-}\textgreater{} L["TIDAK SIGNIFIKAN\textless{}br/\textgreater{}(Hanya Kebetulan)"]}
\NormalTok{    J {-}{-} P{-}Value \textless{} 0.05 {-}{-}\textgreater{} M["SIGNIFIKAN\textless{}br/\textgreater{}(Hubungan Terbukti)"]}
    
\NormalTok{    M {-}{-}\textgreater{} N["Cek Kekuatan Hubungan\textless{}br/\textgreater{}(Odds Ratio / Risk Estimate)"]}
    
\NormalTok{    N {-}{-}\textgreater{} O1["Odds Ratio = 1\textless{}br/\textgreater{}Tidak ngefek"]}
\NormalTok{    N {-}{-}\textgreater{} O2["Odds Ratio \textgreater{} 1\textless{}br/\textgreater{}Risiko/Bahaya Naik"]}
    
\NormalTok{    L {-}{-}\textgreater{} Z(["Selesai"])}
\NormalTok{    O1 {-}{-}\textgreater{} Z}
\NormalTok{    O2 {-}{-}\textgreater{} Z}
\NormalTok{    K {-}{-}\textgreater{} J}
    
\NormalTok{    classDef warning fill:\#ffcccb,stroke:\#ff0000,stroke{-}width:2px;}
\NormalTok{    classDef success fill:\#d4edda,stroke:\#28a745,stroke{-}width:2px;}
\NormalTok{    classDef process fill:\#e1f5fe,stroke:\#0288d1,stroke{-}width:2px;}
    
\NormalTok{    class B,F,J process;}
\NormalTok{    class G,L warning;}
\NormalTok{    class I,M success;}
\end{Highlighting}
\end{Shaded}

\subsubsection{C. Formulasi Matematis: Kalkulasi Konsentrasi Spasial \&
Uji
Chi-Square}\label{c.-formulasi-matematis-kalkulasi-konsentrasi-spasial-uji-chi-square}

Parameterisasi konsentrasi spasial dan pembuktian statistik dihitung
menggunakan sistem formulasi matematis berikut:

\textbf{Persamaan Akumulasi Kapasitas PLTU Kumulatif per Wilayah (MW):}

\begin{Shaded}
\begin{Highlighting}[]
\NormalTok{Kapasitas\_PLTU\_Kumulatif\_t (MW) = Jumlah Kapasitas Aktif Baru (MW) dari Tahun 2014 hingga Tahun t}
\end{Highlighting}
\end{Shaded}

\emph{Keterangan Variabel:} -
\texttt{Kapasitas\_PLTU\_Kumulatif\_t\ (MW)}: Total akumulasi kapasitas
daya terpasang operasional PLTU captive batubara aktif hingga tahun t
(satuan: Megawatt / MW). - \texttt{Kapasitas\ Aktif\ Baru}: Besaran daya
listrik unit PLTU off-grid yang mulai beroperasi komersial pada tahun
tertentu (satuan: Megawatt / MW).

\textbf{Persamaan Rasio Konsentrasi Spasial Fasilitas Smelter (\% pada
Grafik Dashboard):}

\begin{Shaded}
\begin{Highlighting}[]
\NormalTok{Porsi\_Smelter\_Provinsi (\%) = ( Jumlah\_Smelter\_Provinsi / Total\_Smelter\_Sulawesi ) * 100}
\end{Highlighting}
\end{Shaded}

\emph{Keterangan Variabel:} - \texttt{Porsi\_Smelter\_Provinsi\ (\%)}:
Persentase pangsa fasilitas smelter di provinsi bersangkutan terhadap
seluruh Pulau Sulawesi (satuan: Persen / \%). -
\texttt{Jumlah\_Smelter\_Provinsi}: Banyaknya unit smelter yang
beroperasi di wilayah provinsi tertentu. -
\texttt{Total\_Smelter\_Sulawesi}: Total keseluruhan fasilitas smelter
di Pulau Sulawesi (778 unit).

\textbf{Persamaan Uji Independensi Chi-Square Pearson (χ² Kontinjensi
2x2):}

\begin{Shaded}
\begin{Highlighting}[]
\NormalTok{Chi\_Square (χ²) = Jumlah [ (Frekuensi\_Observasi {-} Frekuensi\_Harapan)\^{}2 / Frekuensi\_Harapan ]}
\end{Highlighting}
\end{Shaded}

\emph{Keterangan Variabel:} - \texttt{Chi\_Square\ (χ²)}: Nilai
statistik uji kecocokan Pearson untuk membuktikan ada tidaknya hubungan
ketergantungan antara ekspansi PLTU Captive dengan lonjakan deforestasi
pada panel spasiotemporal (N=60). - \texttt{Frekuensi\_Observasi\ (O)}:
Jumlah kasus aktual yang tercatat pada sel tabel kontinjensi 2x2. -
\texttt{Frekuensi\_Harapan\ (E)}: Jumlah kasus teoretis jika kedua
variabel saling independen: E = (Total Baris * Total Kolom) / N.

\textbf{Persamaan Rasio Keunggulan Risiko (Risk Odds Ratio / OR):}

\begin{Shaded}
\begin{Highlighting}[]
\NormalTok{Odds\_Ratio (OR) = ( a * d ) / ( b * c )}
\end{Highlighting}
\end{Shaded}

\emph{Keterangan Variabel:} - \texttt{Odds\_Ratio\ (OR)}: Ukuran
kelipatan risiko peluang terjadinya deforestasi komoditas tinggi pada
kelompok dengan PLTU Captive aktif (\textgreater0 MW) dibanding kelompok
tanpa PLTU Captive (≤0 MW). - \texttt{a}: Jumlah observasi panel pada
kelompok PLTU Rendah dan Deforestasi Rendah (27 kasus). - \texttt{b}:
Jumlah observasi panel pada kelompok PLTU Rendah dan Deforestasi Tinggi
(10 kasus). - \texttt{c}: Jumlah observasi panel pada kelompok PLTU
Tinggi dan Deforestasi Rendah (3 kasus). - \texttt{d}: Jumlah observasi
panel pada kelompok PLTU Tinggi dan Deforestasi Tinggi (20 kasus).

\subsubsection{D. Matriks Hasil Uji Empiris: Konsentrasi Spasial \&
Skenario
Crosstab}\label{d.-matriks-hasil-uji-empiris-konsentrasi-spasial-skenario-crosstab}

Penerapan sistem pengujian statistik tabulasi silang pada data panel 6
provinsi selama 1 dekade (2014--2023, total 60 observasi) disajikan
secara lengkap pada \textbf{Tabel 1.5} berikut:

\subparagraph{Tabel 1.5: Matriks Tabulasi Silang 2×2, Uji Chi-Square
(χ²), dan Estimasi Odds Ratio Panel PLTU Captive vs Deforestasi
Komoditas
(2014--2023)}\label{tabel-1.5-matriks-tabulasi-silang-22-uji-chi-square-ux3c7uxb2-dan-estimasi-odds-ratio-panel-pltu-captive-vs-deforestasi-komoditas-20142023}

{\def\LTcaptype{none} % do not increment counter
\begin{longtable}[]{@{}
  >{\raggedright\arraybackslash}p{(\linewidth - 12\tabcolsep) * \real{0.1250}}
  >{\centering\arraybackslash}p{(\linewidth - 12\tabcolsep) * \real{0.1562}}
  >{\centering\arraybackslash}p{(\linewidth - 12\tabcolsep) * \real{0.1562}}
  >{\centering\arraybackslash}p{(\linewidth - 12\tabcolsep) * \real{0.1562}}
  >{\raggedright\arraybackslash}p{(\linewidth - 12\tabcolsep) * \real{0.1250}}
  >{\centering\arraybackslash}p{(\linewidth - 12\tabcolsep) * \real{0.1562}}
  >{\raggedright\arraybackslash}p{(\linewidth - 12\tabcolsep) * \real{0.1250}}@{}}
\toprule\noalign{}
\begin{minipage}[b]{\linewidth}\raggedright
Kategori Kapasitas PLTU (X)
\end{minipage} & \begin{minipage}[b]{\linewidth}\centering
Deforestasi Rendah (\textless10,961 Ha)
\end{minipage} & \begin{minipage}[b]{\linewidth}\centering
Deforestasi Tinggi (≥10,961 Ha)
\end{minipage} & \begin{minipage}[b]{\linewidth}\centering
Total Kasus
\end{minipage} & \begin{minipage}[b]{\linewidth}\raggedright
Parameter Statistik Uji
\end{minipage} & \begin{minipage}[b]{\linewidth}\centering
Nilai / df
\end{minipage} & \begin{minipage}[b]{\linewidth}\raggedright
Signifikansi / Kesimpulan
\end{minipage} \\
\midrule\noalign{}
\endhead
\bottomrule\noalign{}
\endlastfoot
\textbf{Rendah (≤0 MW)} & 27 {[}Exp: 18.5{]} & 10 {[}Exp: 18.5{]} & 37
(100\%) & \textbf{Pearson Chi-Square (χ²)} & \textbf{18.049} (df=1) & p
\textless{} 0.0001 (Signifikan) \\
\textbf{Tinggi (\textgreater0 MW)} & 3 {[}Exp: 11.5{]} & 20 {[}Exp:
11.5{]} & 23 (100\%) & \textbf{Likelihood Ratio} & \textbf{19.420}
(df=1) & p \textless{} 0.0001 (Signifikan) \\
\textbf{Total Observasi Panel} & \textbf{30} {[}Exp: 30.0{]} &
\textbf{30} {[}Exp: 30.0{]} & \textbf{60} (100\%) &
\textbf{Linear-by-Linear Association} & \textbf{20.036} (df=1) & p
\textless{} 0.0001 (Signifikan) \\
\textbf{Ukuran Risiko (Risk Estimate)} & Cross-Product: (27×20)/(10×3) &
Rasio Peluang Risiko & OR = 18.00 & \textbf{Odds Ratio (OR)} &
\textbf{18.00x} & \textbf{Risiko Lonjakan 18x Lipat} \\
\end{longtable}
}

\subsubsection{E. Interpretasi Spasial Industri: Eksternalitas dan Efek
Meluber
(Spillover)}\label{e.-interpretasi-spasial-industri-eksternalitas-dan-efek-meluber-spillover}

Hasil pengujian empiris pada Tabel 1.5 membuktikan secara meyakinkan
keterkaitan langsung antara ekspansi PLTU Captive dan kerusakan tutupan
hutan di Pulau Sulawesi:

\begin{enumerate}
\def\labelenumi{\arabic{enumi}.}
\tightlist
\item
  \textbf{Pemusatan Ekstrem di 3 Sentra Ekstraktif Utama:} 100\%
  kapasitas PLTU Captive dan mayoritas smelter berpusat di wilayah ini,
  memicu akumulasi deforestasi komoditas hingga ratusan ribu hektar,
  berbanding terbalik dengan ``Area Non-Smelter''.
\item
  \textbf{Signifikansi Statistik yang Sangat Kuat (p \textless{}
  0.0001):} Hipotesis Nol (H0) ditolak mutlak. Bukti empiris
  mengonfirmasi bahwa penambahan kapasitas PLTU Captive berkorelasi
  langsung dengan lonjakan kehilangan tutupan hutan.
\item
  \textbf{Kelipatan Risiko Bencana Ekologis (Odds Ratio = 18.00x):}
  Wilayah dengan PLTU Captive memiliki risiko deforestasi komoditas 18
  KALI LIPAT lebih besar. Hal ini didorong konversi masif untuk
  infrastruktur pendukung (coal yard, jalur transmisi, dan jalan
  logistik).
\item
  \textbf{Efek Meluber Lintas Batas (Spillover Effect) \& Emisi Karbon
  Terkunci:} Eksternalitas destruktif proyek merambat luas mendegradasi
  DAS dan laut, mengorbankan ruang hidup lokal, serta mengunci emisi
  dari ketergantungan puluhan juta ton batu bara per tahun.
\end{enumerate}

\begin{center}\rule{0.5\linewidth}{0.5pt}\end{center}

\subsection{1.3 Tren Pertumbuhan Izin Tambang Baru \& Uji Signifikansi
Statistik}\label{tren-pertumbuhan-izin-tambang-baru-uji-signifikansi-statistik}

\paragraph{A. Pengantar \& Kerangka
Narasi}\label{a.-pengantar-kerangka-narasi}

Pola perizinan pertambangan di Pulau Sulawesi selama satu dekade
terakhir menunjukkan peningkatan alokasi ruang yang signifikan.
Berdasarkan data agregat \textbf{Minerbaone}, tercatat 574 Izin Usaha
Pertambangan (IUP) baru sepanjang 2014-2024, dengan total luas konsesi
mencapai 819,452 Hektar.

Berdasarkan analisis tren time-series pada grafik ``Penerbitan Izin
Tambang'', penerbitan izin pada periode awal (2014) tercatat lebih
rendah. Peningkatan signifikan terjadi pada periode 2022-2024. Anotasi
pada grafik mencatat kenaikan sebesar \textbf{246\% pada periode
2022-2024}. Data ini mengindikasikan perlunya evaluasi terhadap
instrumen pengendalian perizinan dan tata ruang. Distribusi perizinan
tertinggi berada di Sulawesi Tengah dan Sulawesi Tenggara, yang selaras
dengan kawasan pengembangan industri pemurnian nikel.

Uji \textbf{Crosstabulation} pada analisis ini digunakan untuk mengukur
hubungan antara laju penerbitan perizinan (X) dan indikator deforestasi
di wilayah tersebut (Y).

\paragraph{B. Alur Logika Metodologis
(Flowchart)}\label{b.-alur-logika-metodologis-flowchart}

Pendekatan statistik Time-Series untuk mengidentifikasi tren pertumbuhan
izin tambang diilustrasikan pada \textbf{Bagan Alur 1.5} berikut. Adapun
untuk tahapan analisis inferensial (Uji Chi-Square), alur logikanya
merujuk secara penuh pada \textbf{Bagan Alur 1.4} (di sub-bab
sebelumnya) dengan penyesuaian konfigurasi variabel spesifik sesuai
Tabel Asumsi Dasar di bawah gambar.

Persamaan Formulasi Kurva Parametrik Alur Pelayaran Maritim:

Kurva(t) = (1 - t)\^{}2 * Titik\_Asal + 2 * (1 - t) * t * Titik\_Kontrol
+ t\^{}2 * Titik\_Tujuan, t dalam rentang {[}0, 1{]}

Keterangan Variabel:

• Kurva(t): Vektor posisi koordinat geografis lintasan kapal pada
parameter waktu t.

• Titik\_Asal: Titik koordinat geografis pelabuhan muat khusus di
pesisir Sulawesi.

• Titik\_Kontrol: Titik koordinat jangkar pemandu kurva lengkung di
perairan internasional.

• Titik\_Tujuan: Titik koordinat geografis pelabuhan bongkar di negara
tujuan ekspor.

• t: Parameter interpolasi waktu dan pergerakan lintasan dalam rentang
kontinu {[}0, 1{]}.

SUB-BAB 1.6: Matriks Indikator dan Sumber Data Resmi Bab 1

Seluruh variabel kuantitatif, kategori analisis, satuan ukur, periode
tahun observasi, dan institusi penyedia data primer resmi yang digunakan
dalam Bab 1 dikompilasikan pada Tabel 1.8 berikut:

Tabel 1.8: Matriks Indikator dan Sumber Data Primer Resmi Bab 1
(Ekspansi Industri Ekstraktif)

No

Nama Indikator

Kategori Analisis

Satuan Ukur

Cakupan Tahun

Institusi \& Sumber Data Resmi

1

Izin Usaha Pertambangan (IUP) Baru

Faktor Tekanan Ekstraktif

Unit Izin

2014--2024

Kementerian ESDM (Minerbaone)

2

Luas Wilayah Konsesi Tambang Baru

Faktor Tekanan Ekstraktif

Hektar (Ha)

2014--2024

Kementerian ESDM (Minerbaone)

3

Kapasitas Terpasang PLTU Captive

Infrastruktur Energi Khusus

Megawatt (MW)

2014--2024

Global Energy Monitor (GEM)

4

Fasilitas Pengolahan \& Pemurnian (Smelter)

Fasilitas Industri Hilir

Unit Fasilitas

2014--2024

Kementerian ESDM \& Studi Industri

5

Realisasi Investasi PMDN

Arus Modal Domestik

Triliun Rupiah

2016--2024

Kementerian Investasi / BKPM

6

PDRB Menurut 17 Lapangan Usaha

Struktur Ekonomi Makro

Triliun Rupiah

2016--2024

Badan Pusat Statistik (BPS)

7

PDRB Kabupaten Sentra Tambang

Struktur Ekonomi Daerah

Triliun Rupiah

2016--2024

BPS Kabupaten se-Sulteng

8

Luas Kehilangan Hutan Komoditas

Dampak Tutupan Lahan

Hektar (Ha)

2014--2023

Global Forest Watch (GFW)

9

Simpul Dermaga \& Terminal Khusus Ekspor

Infrastruktur Rantai Pasok

Titik Koordinat \& DWT

2014--2024

KNKT, Perpres PSN, Lap. Terbuka

SUB-BAB 1.7: Bagan Alur Kerangka Kerja Riset Bab 1

Keseluruhan struktur metodologis riset Bab 1 dioperasionalkan melalui
empat fase kerja berurutan sebagaimana disajikan pada Tabel 1.9 berikut:

Tabel 1.9: Matriks Tahapan dan Alur Kerangka Kerja Riset Bab 1

Tahapan Riset

Fokus Metodologis

Bahan \& Sumber Data

Keluaran / Hasil Analisis

Fase I: Pengumpulan Data

Kurasi data resmi lintas kementerian dan lembaga

Publikasi BPS, Minerbaone, BKPM, GEM, dan GFW

Basis Data Tabular Panel Provinsi (2014--2024)

Fase II: Reklasifikasi Hukum

Penyusunan kerangka rantai pasok hukum terintegrasi

UU No.~3/2020, PP No.~96/2021, Perpres No.~112/2022

3 Klaster Makro (Ekstraktif, Akar Rumput, Jasa)

Fase III: Pengujian Statistik

Uji signifikansi hubungan dan rasio peluang

Tabel Kontinjensi, Uji Chi-Square, Odds Ratio

Bukti Kausalitas Signifikan Tekanan vs Deforestasi

Fase IV: Pemetaan Rantai Pasok

Triangulasi data logistik dan pemodelan maritim

Laporan KNKT, Perpres PSN, Kurva Parametrik Bézier

Peta Alur Rantai Pasok Ekspor \& Konsentrasi Spasial 78\%

\subparagraph{Bagan Alur 1.5: Alur Logika Tren Pertumbuhan Izin Tambang
Baru}\label{bagan-alur-1.5-alur-logika-tren-pertumbuhan-izin-tambang-baru}

\begin{Shaded}
\begin{Highlighting}[]
\NormalTok{flowchart TD}
\NormalTok{    subgraph Data\_Preparation["1. Akuisisi \& Penyiapan Data Panel"]}
\NormalTok{        A["Data Penerbitan Izin Baru Minerbaone\textless{}br/\textgreater{}\textless{}i\textgreater{}(Provinsi{-}Tahun, Jumlah Izin, Luas Konsesi)\textless{}/i\textgreater{}"] {-}{-}\textgreater{} C}
\NormalTok{        B["Data Laju Deforestasi GFW\textless{}br/\textgreater{}\textless{}i\textgreater{}(Provinsi{-}Tahun, Deforestasi Komoditas Tambang/Sawit)\textless{}/i\textgreater{}"] {-}{-}\textgreater{} C}
\NormalTok{        C["\textless{}b\textgreater{}Gabung (Merge) Data Panel\textless{}/b\textgreater{}\textless{}br/\textgreater{}6 Provinsi x 10 Tahun = 60 Sampel (Unit Observasi)"]}
\NormalTok{    end}

\NormalTok{    subgraph Time\_Series\_Analysis["2. Analisis Tren (Time{-}Series)"]}
\NormalTok{        C {-}{-}\textgreater{} D["Agregasi Jumlah Izin \& Luas per Tahun"]}
\NormalTok{        D {-}{-}\textgreater{} E["Kalkulasi Pertumbuhan Laju (YoY)"]}
\NormalTok{        E {-}{-}\textgreater{} F["Identifikasi Lonjakan Ekstraktif (Tahun 2022{-}2024)"]}
\NormalTok{    end}

\NormalTok{    subgraph Crosstab\_Analysis["3. Analisis Inferensial (Chi{-}Square)"]}
\NormalTok{        C {-}{-}\textgreater{} G["\textless{}b\textgreater{}Binning Kategori (Threshold Median)\textless{}/b\textgreater{}\textless{}br/\textgreater{}Klasifikasi Tinggi (≥Median) vs Rendah (\textless{}Median)"]}
\NormalTok{        G {-}{-}\textgreater{} H["Definisikan X (Tekanan Ekspansi):\textless{}br/\textgreater{}{-} Jumlah Izin Baru (IUP)\textless{}br/\textgreater{}{-} Luas Konsesi Baru (Ha)"]}
\NormalTok{        H {-}{-}\textgreater{} I["Definisikan Y (Dampak Ekologis):\textless{}br/\textgreater{}{-} Total Deforestasi Alam\textless{}br/\textgreater{}{-} Deforestasi Komoditas"]}
\NormalTok{        I {-}{-}\textgreater{} J["Uji Chi{-}Square (Crosstabulation 2x2)"]}
\NormalTok{    end}

\NormalTok{    subgraph Output["4. Hasil Pengujian Hipotesis"]}
\NormalTok{        F {-}{-}\textgreater{} K["Tren Visual Distribusi Spasial (Bar/Line Chart)"]}
\NormalTok{        J {-}{-}\textgreater{} L\{"P{-}Value \textless{} 0.05?"\}}
\NormalTok{        L {-}{-} YA {-}{-}\textgreater{} M["\textless{}b\textgreater{}SIGNIFIKAN (Tolak H0)\textless{}/b\textgreater{}\textless{}br/\textgreater{}Ada Hubungan: Tekanan Ekspansi Terbukti Meningkatkan Laju Deforestasi"]}
\NormalTok{        L {-}{-} TIDAK {-}{-}\textgreater{} N["\textless{}b\textgreater{}TIDAK SIGNIFIKAN (Gagal Tolak H0)\textless{}/b\textgreater{}\textless{}br/\textgreater{}Indikasi Efek Spillover/Kehancuran Merata"]}
\NormalTok{        M {-}{-}\textgreater{} O["Hitung Odds Ratio (Risk Estimate)"]}
\NormalTok{    end}
\end{Highlighting}
\end{Shaded}

\subparagraph{Tabel 1.5b: Konfigurasi Variabel Uji Chi-Square (Sub-bab
1.3)}\label{tabel-1.5b-konfigurasi-variabel-uji-chi-square-sub-bab-1.3}

{\def\LTcaptype{none} % do not increment counter
\begin{longtable}[]{@{}
  >{\raggedright\arraybackslash}p{(\linewidth - 2\tabcolsep) * \real{0.5000}}
  >{\raggedright\arraybackslash}p{(\linewidth - 2\tabcolsep) * \real{0.5000}}@{}}
\toprule\noalign{}
\begin{minipage}[b]{\linewidth}\raggedright
Komponen Uji
\end{minipage} & \begin{minipage}[b]{\linewidth}\raggedright
Definisi Variabel (Sub-bab 1.3)
\end{minipage} \\
\midrule\noalign{}
\endhead
\bottomrule\noalign{}
\endlastfoot
Variabel Independen (X) & Frekuensi Penerbitan Izin Tambang Baru (IUP) /
Luas Konsesi Baru (Ha) \\
Variabel Dependen (Y) & Deforestasi Komoditas (Ha) / Total Deforestasi
Alam (Ha) \\
Hipotesis Nol (H0) & Tingkat penerbitan izin/luas konsesi tidak
berhubungan dengan laju deforestasi. \\
Hipotesis Alternatif (H1) & Ada hubungan positif antara tingginya
penerbitan izin dengan tingginya laju deforestasi. \\
Threshold Kategori & Nilai Median Data Panel (N=60) \\
\end{longtable}
}

\paragraph{C. Formulasi Matematis: Analisis Tren \& Uji
Chi-Square}\label{c.-formulasi-matematis-analisis-tren-uji-chi-square}

Parameterisasi laju pertumbuhan perizinan dan pengujian signifikansi
dampaknya terhadap deforestasi dihitung menggunakan formulasi berikut:

\textbf{Laju Pertumbuhan Izin Tahunan (Regresi Komparatif YoY):}

\begin{Shaded}
\begin{Highlighting}[]
\NormalTok{Pertumbuhan\_Izin (\%) = [ ( Jumlah\_Izin\_t {-} Jumlah\_Izin\_\{t{-}1\} ) / Jumlah\_Izin\_\{t{-}1\} ] * 100}
\end{Highlighting}
\end{Shaded}

\emph{Keterangan Variabel:} - \texttt{Pertumbuhan\_Izin\ (\%)}:
Persentase perubahan laju penerbitan izin tambang baru antar-tahun
(satuan: Persen / \%). - \texttt{Jumlah\_Izin\_t}: Agregasi jumlah izin
(atau luasan) pada tahun berjalan (t). - \texttt{Jumlah\_Izin\_\{t-1\}}:
Agregasi jumlah izin (atau luasan) pada satu tahun sebelumnya (t - 1).

\textbf{Pengklasifikasian Kategori Data (Binning Threshold Median):}

\begin{Shaded}
\begin{Highlighting}[]
\NormalTok{Kategori = IF(Nilai\_Prov\_Tahun \textgreater{}= Median(Seluruh Panel), "Tinggi", "Rendah")}
\end{Highlighting}
\end{Shaded}

\emph{Keterangan Variabel:} - \texttt{Kategori}: Data panel
spasial-temporal diubah menjadi dua tingkatan untuk uji tabulasi silang
(Tinggi vs Rendah).

Dinamika historis perizinan secara terperinci dapat dilihat pada
\textbf{Tabel 1.5c}, yang menunjukkan tren penerbitan izin baru di
wilayah studi:

\subparagraph{Tabel 1.5c: Tren Penerbitan Izin Tambang Sulawesi
(2014-2024)}\label{tabel-1.5c-tren-penerbitan-izin-tambang-sulawesi-2014-2024}

{\def\LTcaptype{none} % do not increment counter
\begin{longtable}[]{@{}
  >{\centering\arraybackslash}p{(\linewidth - 12\tabcolsep) * \real{0.1429}}
  >{\centering\arraybackslash}p{(\linewidth - 12\tabcolsep) * \real{0.1429}}
  >{\centering\arraybackslash}p{(\linewidth - 12\tabcolsep) * \real{0.1429}}
  >{\centering\arraybackslash}p{(\linewidth - 12\tabcolsep) * \real{0.1429}}
  >{\centering\arraybackslash}p{(\linewidth - 12\tabcolsep) * \real{0.1429}}
  >{\centering\arraybackslash}p{(\linewidth - 12\tabcolsep) * \real{0.1429}}
  >{\centering\arraybackslash}p{(\linewidth - 12\tabcolsep) * \real{0.1429}}@{}}
\toprule\noalign{}
\begin{minipage}[b]{\linewidth}\centering
Tahun
\end{minipage} & \begin{minipage}[b]{\linewidth}\centering
Gorontalo
\end{minipage} & \begin{minipage}[b]{\linewidth}\centering
Sulawesi Barat
\end{minipage} & \begin{minipage}[b]{\linewidth}\centering
Sulawesi Selatan
\end{minipage} & \begin{minipage}[b]{\linewidth}\centering
Sulawesi Tengah
\end{minipage} & \begin{minipage}[b]{\linewidth}\centering
Sulawesi Tenggara
\end{minipage} & \begin{minipage}[b]{\linewidth}\centering
Sulawesi Utara
\end{minipage} \\
\midrule\noalign{}
\endhead
\bottomrule\noalign{}
\endlastfoot
2014 & 0 & 0 & 1 & 6 & 18 & 1 \\
2015 & 0 & 0 & 0 & 3 & 1 & 1 \\
2016 & 0 & 0 & 0 & 5 & 4 & 0 \\
2017 & 0 & 0 & 5 & 7 & 13 & 1 \\
2018 & 1 & 0 & 5 & 7 & 10 & 0 \\
2019 & 0 & 0 & 3 & 2 & 10 & 2 \\
2020 & 0 & 0 & 0 & 12 & 13 & 3 \\
2021 & 1 & 2 & 8 & 17 & 13 & 0 \\
2022 & 1 & 3 & 10 & 31 & 11 & 0 \\
2023 & 1 & 6 & 19 & 83 & 38 & 2 \\
2024 & 3 & 16 & 54 & 87 & 29 & 5 \\
\end{longtable}
}

\paragraph{D. Matriks Hasil Uji Empiris: Ringkasan Skenario
Crosstab}\label{d.-matriks-hasil-uji-empiris-ringkasan-skenario-crosstab}

Hasil lengkap pengujian independensi statistik Chi-Square dan estimasi
Odds Ratio (OR) untuk seluruh faktor tekanan terhadap kehilangan tutupan
hutan komoditas dirangkum pada \textbf{Tabel 1.6} berikut:

\subparagraph{Tabel 1.6: Ringkasan Hasil Uji Independensi Chi-Square
(χ²) dan Odds Ratio (OR) Data Panel Bab
1}\label{tabel-1.6-ringkasan-hasil-uji-independensi-chi-square-ux3c7uxb2-dan-odds-ratio-or-data-panel-bab-1}

{\def\LTcaptype{none} % do not increment counter
\begin{longtable}[]{@{}
  >{\raggedright\arraybackslash}p{(\linewidth - 12\tabcolsep) * \real{0.1250}}
  >{\raggedright\arraybackslash}p{(\linewidth - 12\tabcolsep) * \real{0.1250}}
  >{\centering\arraybackslash}p{(\linewidth - 12\tabcolsep) * \real{0.1562}}
  >{\centering\arraybackslash}p{(\linewidth - 12\tabcolsep) * \real{0.1562}}
  >{\centering\arraybackslash}p{(\linewidth - 12\tabcolsep) * \real{0.1562}}
  >{\centering\arraybackslash}p{(\linewidth - 12\tabcolsep) * \real{0.1562}}
  >{\raggedright\arraybackslash}p{(\linewidth - 12\tabcolsep) * \real{0.1250}}@{}}
\toprule\noalign{}
\begin{minipage}[b]{\linewidth}\raggedright
Variabel Faktor Tekanan
\end{minipage} & \begin{minipage}[b]{\linewidth}\raggedright
Variabel Dampak Lingkungan
\end{minipage} & \begin{minipage}[b]{\linewidth}\centering
Nilai Chi-Square (χ²)
\end{minipage} & \begin{minipage}[b]{\linewidth}\centering
Nilai Signifikansi (p)
\end{minipage} & \begin{minipage}[b]{\linewidth}\centering
Odds Ratio (OR)
\end{minipage} & \begin{minipage}[b]{\linewidth}\centering
Derajat Bebas (df)
\end{minipage} & \begin{minipage}[b]{\linewidth}\raggedright
Kesimpulan Ilmiah
\end{minipage} \\
\midrule\noalign{}
\endhead
\bottomrule\noalign{}
\endlastfoot
Jumlah Izin Tambang Baru (IUP) & Total Deforestasi Alam (Ha) & 17.239 &
\textless{} 0.0001 & 13.75 & 1 & SIGNIFIKAN \\
Jumlah Izin Tambang Baru (IUP) & Deforestasi Komoditas (Ha) & 21.818 &
\textless{} 0.0001 & 21.36 & 1 & SIGNIFIKAN \\
Luas Konsesi Tambang Baru (Ha) & Total Deforestasi Alam (Ha) & 19.267 &
\textless{} 0.0001 & 16.00 & 1 & SIGNIFIKAN \\
Luas Konsesi Tambang Baru (Ha) & Deforestasi Komoditas (Ha) & 19.267 &
\textless{} 0.0001 & 16.00 & 1 & SIGNIFIKAN \\
\end{longtable}
}

\paragraph{E. Analisis Temuan Empiris: Pembedahan Realitas
Ekologis}\label{e.-analisis-temuan-empiris-pembedahan-realitas-ekologis}

Data panel membedah realitas di lapangan: lonjakan izin di wilayah pusat
ekstraksi sejalan dengan tingginya nilai Chi-Square. Nilai Odds Ratio
menegaskan bahwa wilayah dengan tren izin tambang yang tinggi memiliki
peluang lebih besar untuk mengalami tekanan deforestasi tinggi pada
tahun-tahun berjalan dan berikutnya.

Secara spesifik, terjadi \textbf{lonjakan absolut sebesar 246\%} dalam
penerbitan izin tambang baru pada rentang 2022 hingga 2024. Lonjakan
ekstrem ini mengindikasikan percepatan luar biasa dari ekspansi industri
ekstraktif yang mengabaikan kapasitas daya dukung lingkungan tapak,
terutama di sentra-sentra produksi.

\begin{center}\rule{0.5\linewidth}{0.5pt}\end{center}

\subsection{1.4 Analisis Realisasi Investasi PMDN dan Dampak Terhadap
Tutupan
Hutan}\label{analisis-realisasi-investasi-pmdn-dan-dampak-terhadap-tutupan-hutan}

\paragraph{A. Pengantar \& Kerangka
Narasi}\label{a.-pengantar-kerangka-narasi-1}

Akumulasi Penanaman Modal Dalam Negeri sebesar \textbf{Rp 218 Triliun}
(Kementerian Investasi / BKPM) yang masuk dari tahun 2016-2024
berbanding lurus dengan \textbf{1,001,654 Hektar} kehilangan tutupan
hutan komoditas (Global Forest Watch). Grafik sumbu ganda
(\emph{dual-axis}) digunakan untuk membandingkan laju investasi dan laju
deforestasi antara wilayah sentra industri tambang dengan non-sentra.
Terlihat adanya fenomena \textbf{Efek Jeda Waktu (Time-Lagging Effect)},
di mana peningkatan realisasi modal pada tahap awal perizinan dan
konstruksi diikuti oleh lonjakan pembukaan lahan hutan fisik pada 1
hingga 2 tahun berikutnya.

\paragraph{B. Alur Logika Metodologis Analisis Realisasi Investasi
PMDN}\label{b.-alur-logika-metodologis-analisis-realisasi-investasi-pmdn}

Kerangka operasionalisasi uji statistik tabulasi silang antara realisasi
Investasi PMDN dan deforestasi dimodelkan dalam kerangka alur logika
sebagaimana diilustrasikan pada \textbf{Bagan Alur 1.4} berikut:

\subparagraph{Bagan Alur 1.4: Alur Logika Metodologis Uji Independensi
Panel Investasi PMDN vs
Deforestasi}\label{bagan-alur-1.4-alur-logika-metodologis-uji-independensi-panel-investasi-pmdn-vs-deforestasi}

\begin{Shaded}
\begin{Highlighting}[]
\NormalTok{flowchart LR}
\NormalTok{    subgraph Data\_Preparation["1. Akuisisi \& Penyiapan Data Panel"]}
\NormalTok{        A["Data Realisasi Investasi PMDN\textless{}br/\textgreater{}\textless{}i\textgreater{}(Provinsi{-}Tahun, Arus Modal Juta Rp)\textless{}/i\textgreater{}"] {-}{-}\textgreater{} C}
\NormalTok{        B["Data Laju Deforestasi GFW\textless{}br/\textgreater{}\textless{}i\textgreater{}(Provinsi{-}Tahun, Total Deforestasi)\textless{}/i\textgreater{}"] {-}{-}\textgreater{} C}
\NormalTok{        C["\textless{}b\textgreater{}Gabung (Merge) Data Panel\textless{}/b\textgreater{}\textless{}br/\textgreater{}6 Provinsi x 8 Tahun = 48 Sampel Valid"]}
\NormalTok{    end}

\NormalTok{    subgraph Time\_Series\_Analysis["2. Analisis Statistik Tabulasi Silang"]}
\NormalTok{        C {-}{-}\textgreater{} D["Agregasi Median Variabel X \& Y"]}
\NormalTok{        D {-}{-}\textgreater{} E["Kategorisasi Biner\textless{}br/\textgreater{}(Tinggi / Rendah)"]}
\NormalTok{        E {-}{-}\textgreater{} F["Uji Independensi Chi{-}Square Pearson"]}
\NormalTok{        F {-}{-}\textgreater{} G["Kalkulasi Odds Ratio (Risiko Deforestasi)"]}
\NormalTok{    end}
\end{Highlighting}
\end{Shaded}

\paragraph{C. Formulasi Matematis: Agregasi Dampak Ekologis \& Pengujian
Statistik}\label{c.-formulasi-matematis-agregasi-dampak-ekologis-pengujian-statistik}

\textbf{Persamaan Agregasi Luasan Deforestasi Berdasarkan Faktor
Penggerak (Driver)}
\texttt{Total\_Deforestasi\_Driver\_k\ =\ \&sum;(Area\_Loss\_i)\ untuk\ i\ di\ Kategori\_Driver\_k}
- \textbf{Total\_Deforestasi\_Driver\_k:} Total luas kehilangan tutupan
pohon yang diakibatkan oleh faktor penggerak k (contoh: Ekspansi
Komoditas) di seluruh wilayah observasi (satuan: Hektar / Ha). -
\textbf{Kategori\_Driver\_k:} Klasifikasi penyebab utama deforestasi
(Dominant Driver of Tree Cover Loss) berdasarkan model data historis
satelit. - \textbf{Area\_Loss\_i:} Luas kehilangan tutupan pohon pada
piksel observasi ke-i (satuan: Hektar / Ha).

\textbf{Persamaan Perhitungan Akumulasi Kehilangan Hutan Primer (Primary
Forest Loss)}
\texttt{Total\_Primary\_Loss\ =\ \&sum;(Area\_Loss\_j)\ untuk\ j\ di\ mana\ Tipe\_Hutan\ =\ "Primer"}
- \textbf{Total\_Primary\_Loss:} Akumulasi luas konversi tutupan hutan
alam primer tak terganggu (intact primary forest) selama periode
pengamatan (satuan: Hektar / Ha). - \textbf{Tipe\_Hutan:} Klasifikasi
basemap jenis tutupan lahan awal sebelum terjadi deforestasi.

\textbf{Persamaan Estimasi Pelepasan Emisi Karbon (Gross CO2 Emissions)}
\texttt{Emisi\_CO2\_Total\ =\ \&sum;(Area\_Loss\_c\ *\ Faktor\_Emisi\_Biomassa\_c)}
- \textbf{Emisi\_CO2\_Total:} Estimasi agregasi total emisi gas rumah
kaca yang dilepaskan ke atmosfer akibat konversi tutupan (satuan:
Megagrams CO2 / Mg CO2). - \textbf{Faktor\_Emisi\_Biomassa\_c:}
Kandungan karbon rata-rata (above-ground \& below-ground biomass) per
hektar pada koordinat c yang diamati.

Kalkulasi pengujian statistik dihitung menggunakan formulasi Matematis
yang sama dengan Sub-Bab 1.2 dan 1.3, di mana variabel independen (X)
adalah \textbf{Investasi PMDN (Juta Rp)} dan variabel dependen (Y)
adalah \textbf{Deforestasi Komoditas (Hektar)}.

\textbf{Persamaan Kategorisasi Nilai Ambang Batas Median:}

\begin{Shaded}
\begin{Highlighting}[]
\NormalTok{{-} Jika Nilai \textgreater{} Median, maka Kategori = Tinggi}
\NormalTok{{-} Jika Nilai \textless{}= Median, maka Kategori = Rendah}
\end{Highlighting}
\end{Shaded}

\paragraph{D. Matriks Hasil Uji Empiris: Konsentrasi Spasial \& Skenario
Crosstab}\label{d.-matriks-hasil-uji-empiris-konsentrasi-spasial-skenario-crosstab-1}

Tingkat alokasi konsesi dan dampaknya terhadap tutupan hutan dapat
dilihat secara empiris melalui perbandingan luas konsesi baru di Daerah
Sentra Tambang (Morowali \& Konawe) dengan wilayah non-sentra pada
\textbf{Tabel 1.7b} berikut:

\subparagraph{Tabel 1.7b: Representasi Spasial Luas Konsesi Baru (Ha) di
Daerah Sentra Tambang vs Non-Sentra
(2014-2023)}\label{tabel-1.7b-representasi-spasial-luas-konsesi-baru-ha-di-daerah-sentra-tambang-vs-non-sentra-2014-2023}

{\def\LTcaptype{none} % do not increment counter
\begin{longtable}[]{@{}
  >{\centering\arraybackslash}p{(\linewidth - 4\tabcolsep) * \real{0.3333}}
  >{\centering\arraybackslash}p{(\linewidth - 4\tabcolsep) * \real{0.3333}}
  >{\centering\arraybackslash}p{(\linewidth - 4\tabcolsep) * \real{0.3333}}@{}}
\toprule\noalign{}
\begin{minipage}[b]{\linewidth}\centering
Tahun
\end{minipage} & \begin{minipage}[b]{\linewidth}\centering
Luas Konsesi Baru Daerah Sentra Tambang (Ha)
\end{minipage} & \begin{minipage}[b]{\linewidth}\centering
Luas Konsesi Baru Daerah Non-Sentra (Ha)
\end{minipage} \\
\midrule\noalign{}
\endhead
\bottomrule\noalign{}
\endlastfoot
2014 & 39,216.7 & 10,301.4 \\
2015 & 13,370.0 & 8,969.0 \\
2016 & 12,515.8 & 0.0 \\
2017 & 121,742.9 & 57,722.2 \\
2018 & 15,634.7 & 22,336.0 \\
2019 & 32,458.3 & 29,509.9 \\
2020 & 60,420.8 & 46,139.0 \\
2021 & 28,835.0 & 1,588.4 \\
2022 & 58,268.1 & 7,859.4 \\
2023 & 52,138.7 & 16,418.2 \\
\end{longtable}
}

Terkait dengan hilangnya luasan hutan tersebut, pembedahan lebih lanjut
berdasarkan aktor utama, luasan hutan primer, dan estimasi emisi karbon
komoditas dapat dilihat pada \textbf{Tabel 1.7c} berikut:

\subparagraph{Tabel 1.7c: Matriks Pembedahan Ekologis Aktor \& Emisi
Karbon (Periode
2001-2025)}\label{tabel-1.7c-matriks-pembedahan-ekologis-aktor-emisi-karbon-periode-2001-2025}

{\def\LTcaptype{none} % do not increment counter
\begin{longtable}[]{@{}
  >{\raggedright\arraybackslash}p{(\linewidth - 4\tabcolsep) * \real{0.2857}}
  >{\centering\arraybackslash}p{(\linewidth - 4\tabcolsep) * \real{0.3571}}
  >{\centering\arraybackslash}p{(\linewidth - 4\tabcolsep) * \real{0.3571}}@{}}
\toprule\noalign{}
\begin{minipage}[b]{\linewidth}\raggedright
Kategori Aktor / Metrik Ekologis
\end{minipage} & \begin{minipage}[b]{\linewidth}\centering
Nilai Agregat
\end{minipage} & \begin{minipage}[b]{\linewidth}\centering
Persentase dari Total Kehilangan
\end{minipage} \\
\midrule\noalign{}
\endhead
\bottomrule\noalign{}
\endlastfoot
\textbf{Ekspansi Komoditas (Tambang \& Sawit)} & 1,890,659 Hektar &
48.4\% \\
\textbf{Kehutanan (Logging)} & 247,011 Hektar & 6.3\% \\
\textbf{Pertanian Berpindah} & 115,404 Hektar & 2.9\% \\
\textbf{Total Kehilangan Hutan Primer} & \textbf{3,904,079 Hektar} &
\textbf{100.0\%} \\
\textbf{Estimasi Emisi Karbon Komoditas} & 1,282,195,705 Mg CO2 & - \\
\end{longtable}
}

\subparagraph{Tabel 1.7d: Konfigurasi Variabel Uji Chi-Square (Sub-bab
1.4)}\label{tabel-1.7d-konfigurasi-variabel-uji-chi-square-sub-bab-1.4}

{\def\LTcaptype{none} % do not increment counter
\begin{longtable}[]{@{}
  >{\raggedright\arraybackslash}p{(\linewidth - 2\tabcolsep) * \real{0.5000}}
  >{\raggedright\arraybackslash}p{(\linewidth - 2\tabcolsep) * \real{0.5000}}@{}}
\toprule\noalign{}
\begin{minipage}[b]{\linewidth}\raggedright
Komponen Uji
\end{minipage} & \begin{minipage}[b]{\linewidth}\raggedright
Definisi Variabel (Sub-bab 1.4)
\end{minipage} \\
\midrule\noalign{}
\endhead
\bottomrule\noalign{}
\endlastfoot
Variabel Independen (X) & Realisasi Investasi PMDN (Juta Rupiah) \\
Variabel Dependen (Y) & Total Deforestasi Alam (Ha) / Deforestasi
Komoditas (Ha) \\
Hipotesis Nol (H0) & Tingginya realisasi investasi PMDN tidak
berhubungan dengan laju deforestasi. \\
Hipotesis Alternatif (H1) & Ada hubungan positif antara tingginya
realisasi investasi PMDN dengan laju deforestasi. \\
Threshold Kategori & Nilai Median Data Panel (N=48) \\
\end{longtable}
}

Tabel di bawah ini merangkum hasil pengujian statistik (Chi-Square)
untuk semua kemungkinan kombinasi indikator antara Realisasi Investasi
PMDN dan Dampak Ekologis pada panel data 2016-2023. Hasil tersebut
ditampilkan pada \textbf{Tabel 1.8} berikut:

\subparagraph{Tabel 1.8: Ringkasan Eksekutif Seluruh Skenario Crosstab
Realisasi Investasi PMDN Bab
1}\label{tabel-1.8-ringkasan-eksekutif-seluruh-skenario-crosstab-realisasi-investasi-pmdn-bab-1}

{\def\LTcaptype{none} % do not increment counter
\begin{longtable}[]{@{}
  >{\raggedright\arraybackslash}p{(\linewidth - 10\tabcolsep) * \real{0.1481}}
  >{\raggedright\arraybackslash}p{(\linewidth - 10\tabcolsep) * \real{0.1481}}
  >{\centering\arraybackslash}p{(\linewidth - 10\tabcolsep) * \real{0.1852}}
  >{\centering\arraybackslash}p{(\linewidth - 10\tabcolsep) * \real{0.1852}}
  >{\centering\arraybackslash}p{(\linewidth - 10\tabcolsep) * \real{0.1852}}
  >{\raggedright\arraybackslash}p{(\linewidth - 10\tabcolsep) * \real{0.1481}}@{}}
\toprule\noalign{}
\begin{minipage}[b]{\linewidth}\raggedright
Variabel Independen (X)
\end{minipage} & \begin{minipage}[b]{\linewidth}\raggedright
Variabel Dependen (Y)
\end{minipage} & \begin{minipage}[b]{\linewidth}\centering
Chi-Square (χ²)
\end{minipage} & \begin{minipage}[b]{\linewidth}\centering
P-Value
\end{minipage} & \begin{minipage}[b]{\linewidth}\centering
Odds Ratio
\end{minipage} & \begin{minipage}[b]{\linewidth}\raggedright
Kesimpulan
\end{minipage} \\
\midrule\noalign{}
\endhead
\bottomrule\noalign{}
\endlastfoot
Realisasi Investasi PMDN (Juta Rp) & Total Deforestasi Alam (Hektar) &
2.042 & p = 0.1530 & 2.0 & TIDAK SIGNIFIKAN \\
Realisasi Investasi PMDN (Juta Rp) & Deforestasi Komoditas Tambang/Sawit
(Hektar) & 3.375 & p = 0.0662 & 2.3 & TIDAK SIGNIFIKAN \\
\end{longtable}
}

\paragraph{E. Analisis Temuan Empiris: Efek Jeda Waktu
(Time-Lagging)}\label{e.-analisis-temuan-empiris-efek-jeda-waktu-time-lagging}

Hasil pengujian seluruh skenario tabulasi silang PMDN mengungkap
fenomena yang kompleks dalam alur investasi ekstraktif:

\begin{enumerate}
\def\labelenumi{\arabic{enumi}.}
\tightlist
\item
  \textbf{Ketidaksignifikanan Simultan \& Variasi P-Value:} Tingkat
  signifikansi yang bervariasi menyingkap tabir jeda waktu (lagging
  effect) dalam eksekusi investasi di lapangan.
\item
  \textbf{Jeda Waktu Eksekusi Investasi (Lagging Effect):} Suntikan
  modal masif di tahun tertentu tidak secara instan berwujud pembabatan
  lahan di tahun yang sama. Modal tersebut tertahan pada fase birokrasi,
  pembebasan lahan, dan pengadaan infrastruktur, sebelum daya rusaknya
  mengonversi lanskap hutan pada tahun-tahun berikutnya.
\item
  \textbf{Konsentrasi Modal Ekstrem di 3 Provinsi:} Data spasial
  membuktikan 89\% dari total modal PMDN ekstraktif se-Sulawesi hanya
  tersedot ke tiga provinsi sentra (Sulteng, Sultra, Sulsel),
  mengakibatkan polarisasi pertumbuhan dan mengunci ketimpangan spasial.
\end{enumerate}

\begin{center}\rule{0.5\linewidth}{0.5pt}\end{center}

\subsection{1.5 Pelabuhan Ekspor \& Peta Jalur Distribusi Logistik Nikel
Sulawesi}\label{pelabuhan-ekspor-peta-jalur-distribusi-logistik-nikel-sulawesi}

\paragraph{A. Pengantar \& Kerangka
Narasi}\label{a.-pengantar-kerangka-narasi-2}

Eksploitasi nikel di Sulawesi tidak berhenti di tapak darat, melainkan
terhubung langsung ke pasar global melalui infrastruktur pelabuhan.
Bagian ini memetakan simpul logistik maritim yang mendistribusikan
produk ekstraktif (NPI, Matte, MHP) dari pesisir Sulawesi menuju negara
tujuan utama seperti Tiongkok dan Jepang.

\paragraph{B. Alur Logika Metodologis (Validasi Silang \& Kurva
Bézier)}\label{b.-alur-logika-metodologis-validasi-silang-kurva-buxe9zier}

Verifikasi titik pelabuhan dan terminal khusus ekspor nikel dilakukan
melalui protokol triangulasi informasi publik sebagaimana
divisualisasikan pada \textbf{Bagan Alur 1.5} berikut:

\subparagraph{Bagan Alur 1.5: Alur Logika Metodologis Validasi Silang
(OSINT) dan Pemetaan Spasial
Pelabuhan}\label{bagan-alur-1.5-alur-logika-metodologis-validasi-silang-osint-dan-pemetaan-spasial-pelabuhan}

\begin{Shaded}
\begin{Highlighting}[]
\NormalTok{flowchart LR}
\NormalTok{    subgraph Data\_Acquisition["1. Validasi Silang Dokumen Publik (OSINT)"]}
\NormalTok{        A["Laporan KNKT\textless{}br/\textgreater{}\textless{}i\textgreater{}(Kapasitas Muatan \& DWT)\textless{}/i\textgreater{}"] {-}{-}\textgreater{} D}
\NormalTok{        B["Lampiran Perpres 109/2020\textless{}br/\textgreater{}\textless{}i\textgreater{}(Status Proyek Strategis Nasional)\textless{}/i\textgreater{}"] {-}{-}\textgreater{} D}
\NormalTok{        C["Laporan Tahunan \& Laporan Keberlanjutan\textless{}br/\textgreater{}\textless{}i\textgreater{}(Fasilitas Terminal Khusus)\textless{}/i\textgreater{}"] {-}{-}\textgreater{} D}
\NormalTok{    end}

\NormalTok{    subgraph Data\_Processing["2. Triangulasi \& Pemetaan Rute"]}
\NormalTok{        D["\textless{}b\textgreater{}Inventarisasi Simpul Pelabuhan\textless{}/b\textgreater{}\textless{}br/\textgreater{}(Identifikasi 6 Kawasan Utama)"] {-}{-}\textgreater{} E["Ekstraksi Koordinat Geografis\textless{}br/\textgreater{}(Latitude \& Longitude)"]}
\NormalTok{        E {-}{-}\textgreater{} F["Pemodelan Kurva Parametrik (Bézier Curve)"]}
\NormalTok{        F {-}{-}\textgreater{} G["Visualisasi Spatial Logistic Mapping"]}
\NormalTok{    end}
\end{Highlighting}
\end{Shaded}

Detail dari keempat sumber informasi validasi silang (triangulasi)
publik adalah sebagai berikut:

\begin{enumerate}
\def\labelenumi{\arabic{enumi}.}
\tightlist
\item
  \textbf{Laporan Investigasi Keselamatan Transportasi Laut (KNKT):}
  Memverifikasi kapasitas dermaga curah dan bobot muatan kapal
  pengangkut bijih nikel hingga 52.378 DWT.
\item
  \textbf{Regulasi Proyek Strategis Nasional (PSN):} Lampiran Perpres
  No.~109 Tahun 2020 sektor kawasan industri terpadu.
\item
  \textbf{Laporan Keberlanjutan \& Tahunan Korporasi Terbuka:} Laporan
  resmi PT Vale Indonesia Tbk dan PT ANTAM Tbk mengenai fasilitas
  pelabuhan khusus.
\item
  \textbf{Laporan Audit Lembaga Riset Independen:} Publikasi riset
  independen mengenai rantai pasok dan operasional terminal khusus
  maritim.
\end{enumerate}

\paragraph{C. Formulasi Matematis: Kurva Parametrik Alur
Pelayaran}\label{c.-formulasi-matematis-kurva-parametrik-alur-pelayaran}

Pemetaan alur pelayaran internasional dimodelkan menggunakan formulasi
kurva parametrik lengkung (\emph{Bézier Curve}) untuk merepresentasikan
alur laut kepulauan dan rute maritim internasional secara realistis:

\textbf{Persamaan Formulasi Kurva Parametrik Alur Pelayaran Maritim:}

\begin{Shaded}
\begin{Highlighting}[]
\NormalTok{Kurva(t) = (1 {-} t)\^{}2 * Titik\_Asal + 2 * (1 {-} t) * t * Titik\_Kontrol + t\^{}2 * Titik\_Tujuan}
\end{Highlighting}
\end{Shaded}

\emph{Keterangan Variabel:} - \texttt{Kurva(t)}: Vektor posisi koordinat
geografis lintasan kapal pada parameter waktu t (rentang kontinu {[}0,
1{]}). - \texttt{Titik\_Asal}: Titik koordinat geografis pelabuhan muat
khusus di pesisir Sulawesi. - \texttt{Titik\_Kontrol}: Titik koordinat
jangkar pemandu kurva lengkung di perairan internasional. -
\texttt{Titik\_Tujuan}: Titik koordinat geografis pelabuhan bongkar di
negara tujuan ekspor.

\paragraph{D. Matriks Hasil Uji Empiris: Inventarisasi 6 Simpul
Pelabuhan}\label{d.-matriks-hasil-uji-empiris-inventarisasi-6-simpul-pelabuhan}

Berdasarkan protokol validasi silang tersebut, profil komprehensif enam
simpul pelabuhan dan terminal khusus utama di Pulau Sulawesi dipetakan
pada \textbf{Tabel 1.7} berikut:

\subparagraph{Tabel 1.7: Inventarisasi Enam Simpul Pelabuhan dan
Terminal Khusus Ekspor Nikel di Pulau
Sulawesi}\label{tabel-1.7-inventarisasi-enam-simpul-pelabuhan-dan-terminal-khusus-ekspor-nikel-di-pulau-sulawesi}

{\def\LTcaptype{none} % do not increment counter
\begin{longtable}[]{@{}
  >{\raggedright\arraybackslash}p{(\linewidth - 10\tabcolsep) * \real{0.1538}}
  >{\raggedright\arraybackslash}p{(\linewidth - 10\tabcolsep) * \real{0.1538}}
  >{\raggedright\arraybackslash}p{(\linewidth - 10\tabcolsep) * \real{0.1538}}
  >{\centering\arraybackslash}p{(\linewidth - 10\tabcolsep) * \real{0.1923}}
  >{\centering\arraybackslash}p{(\linewidth - 10\tabcolsep) * \real{0.1923}}
  >{\raggedright\arraybackslash}p{(\linewidth - 10\tabcolsep) * \real{0.1538}}@{}}
\toprule\noalign{}
\begin{minipage}[b]{\linewidth}\raggedright
Simpul Kawasan Industri
\end{minipage} & \begin{minipage}[b]{\linewidth}\raggedright
Wilayah Administrasi
\end{minipage} & \begin{minipage}[b]{\linewidth}\raggedright
Fasilitas Pelabuhan / Terminal
\end{minipage} & \begin{minipage}[b]{\linewidth}\centering
Status Regulasi
\end{minipage} & \begin{minipage}[b]{\linewidth}\centering
Kapasitas Kapal
\end{minipage} & \begin{minipage}[b]{\linewidth}\raggedright
Tujuan Utama Ekspor
\end{minipage} \\
\midrule\noalign{}
\endhead
\bottomrule\noalign{}
\endlastfoot
\textbf{IMIP Morowali} & Morowali, Sulawesi Tengah & Pelabuhan Samudera
\& Dermaga Curah & PSN (Perpres 109/2020) & Hingga 52.378 DWT & Pasar
Global (Tiongkok) \\
\textbf{GNI Morowali Utara} & Morowali Utara, Sulteng & Terminal Khusus
Pesisir Tomori & Izin Industri Mandiri & Hingga 30.000 DWT & Pasar
Global (Tiongkok) \\
\textbf{VDNI Konawe} & Konawe, Sulawesi Tenggara & Dermaga Khusus Curah
\& Kargo & PSN (Perpres 109/2020) & Hingga 50.000 DWT & Pasar Global
(Tiongkok) \\
\textbf{OSS Konawe} & Konawe, Sulawesi Tenggara & Dermaga Terintegrasi
Konawe & PSN (Perpres 109/2020) & Hingga 50.000 DWT & Pasar Global
(Tiongkok) \\
\textbf{Pomalaa (ANTAM)} & Kolaka, Sulawesi Tenggara & Dermaga Pomalaa
\& Konveyor & Kawasan BUMN Industri & Hingga 12.000 DWT & Jepang \&
Korsel \\
\textbf{Sorowako (Vale)} & Luwu Timur, Sulawesi Selatan & Pelabuhan
Balantang Malili & Kontrak Karya Tambang & Hingga 15.000 DWT & Jepang \&
Skandinavia \\
\end{longtable}
}

\paragraph{E. Interpretasi Spasial Industri: Rantai Pasok
Global}\label{e.-interpretasi-spasial-industri-rantai-pasok-global}

Visualisasi jalur pelayaran ini mengonfirmasi temuan struktural pada
sub-bab sebelumnya, di mana eksploitasi ekstraktif sepenuhnya
berorientasi ekspor:

\begin{enumerate}
\def\labelenumi{\arabic{enumi}.}
\tightlist
\item
  \textbf{Infrastruktur Raksasa Penunjang Ekspor:} Keenam simpul
  pelabuhan melayani kapal dengan kapasitas hingga 50.000 DWT,
  memastikan arus bahan baku mentah (raw materials) terus mengalir ke
  pusat-pusat industri global, terutama Tiongkok.
\item
  \textbf{Eksternalitas bagi Ekosistem Pesisir:} Pembangunan pelabuhan
  samudera di pesisir Morowali dan Konawe memicu fragmentasi ekosistem
  laut, perusakan terumbu karang, dan konflik ruang dengan nelayan
  tangkap tradisional.
\end{enumerate}

\begin{itemize}
\tightlist
\item
  \texttt{Titik\_Asal}: Titik koordinat geografis pelabuhan muat khusus.
\item
  \texttt{Titik\_Kontrol}: Titik koordinat jangkar pemandu kurva
  lengkung.
\item
  \texttt{Titik\_Tujuan}: Titik koordinat geografis pelabuhan bongkar
  negara tujuan.
\item
  \texttt{t}: Parameter interpolasi waktu (rentang 0-1).
\end{itemize}

\begin{center}\rule{0.5\linewidth}{0.5pt}\end{center}

\subsection{1.6 Matriks Indikator dan Sumber Data Resmi Bab
1}\label{matriks-indikator-dan-sumber-data-resmi-bab-1}

Seluruh variabel kuantitatif, kategori analisis, satuan ukur, periode
tahun observasi, dan institusi penyedia data primer resmi yang digunakan
dalam Bab 1 dikompilasikan pada \textbf{Tabel 1.8} berikut:

\subparagraph{Tabel 1.8: Matriks Indikator dan Sumber Data Primer Resmi
Bab
1}\label{tabel-1.8-matriks-indikator-dan-sumber-data-primer-resmi-bab-1}

{\def\LTcaptype{none} % do not increment counter
\begin{longtable}[]{@{}
  >{\centering\arraybackslash}p{(\linewidth - 10\tabcolsep) * \real{0.1852}}
  >{\raggedright\arraybackslash}p{(\linewidth - 10\tabcolsep) * \real{0.1481}}
  >{\raggedright\arraybackslash}p{(\linewidth - 10\tabcolsep) * \real{0.1481}}
  >{\centering\arraybackslash}p{(\linewidth - 10\tabcolsep) * \real{0.1852}}
  >{\centering\arraybackslash}p{(\linewidth - 10\tabcolsep) * \real{0.1852}}
  >{\raggedright\arraybackslash}p{(\linewidth - 10\tabcolsep) * \real{0.1481}}@{}}
\toprule\noalign{}
\begin{minipage}[b]{\linewidth}\centering
No
\end{minipage} & \begin{minipage}[b]{\linewidth}\raggedright
Nama Indikator
\end{minipage} & \begin{minipage}[b]{\linewidth}\raggedright
Kategori Analisis
\end{minipage} & \begin{minipage}[b]{\linewidth}\centering
Satuan Ukur
\end{minipage} & \begin{minipage}[b]{\linewidth}\centering
Cakupan Tahun
\end{minipage} & \begin{minipage}[b]{\linewidth}\raggedright
Institusi \& Sumber Data Resmi
\end{minipage} \\
\midrule\noalign{}
\endhead
\bottomrule\noalign{}
\endlastfoot
1 & Izin Usaha Pertambangan (IUP) Baru & Faktor Tekanan Ekstraktif &
Unit Izin & 2014--2024 & Kementerian ESDM (Minerbaone) \\
2 & Luas Wilayah Konsesi Tambang Baru & Faktor Tekanan Ekstraktif &
Hektar (Ha) & 2014--2024 & Kementerian ESDM (Minerbaone) \\
3 & Kapasitas Terpasang PLTU Captive & Infrastruktur Energi Khusus &
Megawatt (MW) & 2014--2024 & Global Energy Monitor (GEM) \\
4 & Fasilitas Pengolahan \& Pemurnian (Smelter) & Fasilitas Industri
Hilir & Unit Fasilitas & 2014--2024 & Kementerian ESDM \& Studi
Industri \\
5 & Realisasi Investasi PMDN & Arus Modal Domestik & Triliun Rupiah &
2016--2024 & Kementerian Investasi / BKPM \\
6 & PDRB Menurut 17 Lapangan Usaha & Struktur Ekonomi Makro & Triliun
Rupiah & 2016--2024 & Badan Pusat Statistik (BPS Provinsi) \\
7 & PDRB Kabupaten Sentra Tambang & Struktur Ekonomi Daerah & Triliun
Rupiah & 2016--2024 & BPS Kabupaten se-Sulteng \\
8 & Luas Kehilangan Hutan Komoditas & Dampak Tutupan Lahan & Hektar (Ha)
& 2014--2023 & Global Forest Watch (GFW) \\
9 & Simpul Dermaga \& Terminal Khusus Ekspor & Infrastruktur Rantai
Pasok & Titik Koordinat \& DWT & 2014--2024 & KNKT, Perpres PSN, Lap.
Terbuka \\
\end{longtable}
}

\begin{center}\rule{0.5\linewidth}{0.5pt}\end{center}

\subsection{1.7 Bagan Alur Kerangka Kerja Riset Bab
1}\label{bagan-alur-kerangka-kerja-riset-bab-1}

Keseluruhan struktur metodologis riset Bab 1 dioperasionalkan melalui
empat fase kerja berurutan sebagaimana disajikan pada \textbf{Tabel 1.9}
berikut:

\subparagraph{Tabel 1.9: Matriks Tahapan dan Alur Kerangka Kerja Riset
Bab
1}\label{tabel-1.9-matriks-tahapan-dan-alur-kerangka-kerja-riset-bab-1}

{\def\LTcaptype{none} % do not increment counter
\begin{longtable}[]{@{}
  >{\raggedright\arraybackslash}p{(\linewidth - 6\tabcolsep) * \real{0.2500}}
  >{\raggedright\arraybackslash}p{(\linewidth - 6\tabcolsep) * \real{0.2500}}
  >{\raggedright\arraybackslash}p{(\linewidth - 6\tabcolsep) * \real{0.2500}}
  >{\raggedright\arraybackslash}p{(\linewidth - 6\tabcolsep) * \real{0.2500}}@{}}
\toprule\noalign{}
\begin{minipage}[b]{\linewidth}\raggedright
Tahapan Riset
\end{minipage} & \begin{minipage}[b]{\linewidth}\raggedright
Fokus Metodologis
\end{minipage} & \begin{minipage}[b]{\linewidth}\raggedright
Bahan \& Sumber Data
\end{minipage} & \begin{minipage}[b]{\linewidth}\raggedright
Keluaran / Hasil Analisis
\end{minipage} \\
\midrule\noalign{}
\endhead
\bottomrule\noalign{}
\endlastfoot
\textbf{Fase I: Pengumpulan Data} & Kurasi data resmi lintas kementerian
dan lembaga & Publikasi BPS, Minerbaone, BKPM, GEM, dan GFW & Basis Data
Tabular Panel Provinsi (2014--2024) \\
\textbf{Fase II: Reklasifikasi Hukum} & Penyusunan kerangka rantai pasok
hukum terintegrasi & UU No.~3/2020, PP No.~96/2021, Perpres No.~112/2022
& 3 Klaster Makro (Ekstraktif, Akar Rumput, Jasa) \\
\textbf{Fase III: Pengujian Statistik} & Uji signifikansi hubungan dan
rasio peluang & Tabel Kontinjensi, Uji Chi-Square, Odds Ratio & Bukti
Kausalitas Signifikan Tekanan vs Deforestasi \\
\textbf{Fase IV: Pemetaan Rantai Pasok} & Triangulasi data logistik dan
pemodelan maritim & Laporan KNKT, Perpres PSN, Kurva Parametrik Bézier &
Peta Alur Rantai Pasok Ekspor \& Konsentrasi Spasial 78\% \\
\end{longtable}
}

\end{document}
